\documentclass[11pt,a4paper]{article}
\usepackage[T1]{fontenc}
\usepackage[utf8]{inputenc}
\usepackage[french]{babel}
\usepackage{lmodern}
\usepackage{amsmath,amssymb}
\usepackage[margin=2.35cm]{geometry}
\usepackage{booktabs,longtable,tabularx}
\usepackage{array}
\usepackage{enumitem}
\usepackage{graphicx}
\usepackage{xcolor}
\usepackage{microtype}
\usepackage{float}
\usepackage{hyperref}
% Caractères Unicode présents dans les fragments générés (noms de fenêtres).
\DeclareUnicodeCharacter{2264}{\ensuremath{\le}}
\DeclareUnicodeCharacter{2265}{\ensuremath{\ge}}
\DeclareUnicodeCharacter{03C3}{\ensuremath{\sigma}}
\DeclareUnicodeCharacter{03B3}{\ensuremath{\gamma}}
\DeclareUnicodeCharacter{03C6}{\ensuremath{\varphi}}

\hypersetup{colorlinks=true,linkcolor=blue!45!black,urlcolor=blue!45!black}
\setlist{nosep}
\newcommand{\code}[1]{\texttt{#1}}
\newcommand{\Kaut}{K^{*}_{\mathrm{aut}}}
\newcommand{\fait}{\textbf{[Fait observé]}}
\newcommand{\inference}{\textbf{[Inférence]}}
\newcommand{\hyp}{\textbf{[Hypothèse]}}
\newcommand{\incertitude}{\textbf{[Incertitude]}}

\title{M4.3Live --- rapport de conception\\
\large Institution de principal « production jointe maximale »,
interventions en direct, déterminisme}
\author{Programme M4.3Live --- dossier \code{m4\_3live\_credit\_soc}}
\date{17 août 2026}

\begin{document}
\maketitle

\begin{abstract}
Ce document justifie les décisions d'architecture de M4.3Live et tranche,
explicitement, les huit points que le prompt laissait ouverts (§11). Les
résultats scientifiques sont dans \code{report/rapport\_final.pdf} ; on ne
trouvera ici que ce qui concerne la construction du système.

Trois décisions structurent le reste. \textbf{(1)} L'institution de
principal est routée sur des identifiants \emph{entiers} de technologie,
et le cas « même technologie » applique littéralement $(K_\ell-K_b)/2$ :
la baseline homogène de M4.3Live est donc bit à bit celle de M4.3 par
construction, ce qui a été vérifié sur les 8000 pas et les 26 colonnes
d'un run stocké. \textbf{(2)} La compilation des tables du noyau est
\emph{synchrone}, contre la suggestion du rapport d'architecture de la
déporter sur un thread : c'est la condition pour que le rejeu d'un journal
d'interventions soit strictement identique. \textbf{(3)} Le chemin
« tiède » du rapport est implémenté et testé, mais n'est pas le défaut :
son erreur mesurée sur le domaine réel est de $9{,}8\cdot10^{-2}$ unité de
capital, sept ordres de grandeur au-dessus de la table cubique, pour un
gain de débit que la part du noyau dans le coût d'un pas ne justifie pas.
\end{abstract}

\tableofcontents

% Section partagée par les deux rapports M4.3Live : le modèle et les
% notations, pour que chaque document se lise sans accès au code ni aux
% rapports des lignées antérieures.
\section{Le modèle, et les notations employées}
\label{sec:modele}

\subsection{Ce que simule M4.3Live}

Le modèle décrit une population d'entités économiques abstraites qui
échangent une ressource unique, le \emph{joule}, assimilée ici à un capital.
Il n'y a ni monnaie, ni banque centrale, ni prix : seulement des capitaux,
une fonction de production, et des prêts bilatéraux perpétuels.

Une entité $i$ est entièrement décrite par trois nombres : son capital
$K_i \ge 0$ et sa \emph{technologie} $(A_i, \gamma_i)$, qui fixe la quantité
qu'elle produit à chaque pas de temps,
\begin{equation}
\text{production de } i \;=\; A_i\,K_i^{\gamma_i},
\qquad A_i > 0,\quad 0 < \gamma_i < 1 .
\label{eq:production}
\end{equation}
L'exposant $\gamma_i < 1$ rend les rendements \emph{marginaux} décroissants :
un joule supplémentaire rapporte d'autant moins que le capital est déjà
grand ($\partial^2/\partial K^2 < 0$). La production, elle, continue de
croître avec le capital --- doubler le capital augmente la production, mais
de moins du double. C'est cette concavité qui crée un intérêt à déplacer du
capital d'une entité riche vers une entité pauvre --- le fondement même du
marché du crédit du modèle.

\subsection{Un pas de temps}

Chaque pas déroule huit phases, dans cet ordre invariable :

\begin{enumerate}
\item \textbf{Interventions} --- les changements de paramètres demandés en
  direct sont appliqués (propre à M4.3Live ; voir §\ref{sec:portees} pour
  leurs portées).
\item \textbf{Naissances} --- un nombre $\mathrm{Poisson}(\lambda)$ d'entités
  neuves apparaissent, chacune dotée du capital de naissance $K_0$.
\item \textbf{Choc} --- chaque capital est multiplié par $e^{\xi}$ avec
  $\xi \sim \mathcal{N}(-\sigma^2/2,\ \sigma^2)$, un choc multiplicatif
  d'espérance 1 (il ne crée ni ne détruit de richesse en moyenne).
  \emph{Ordre de grandeur.} L'écart typique se lit sur
  $\mathbb{E}|\xi| = \sigma\sqrt{2/\pi}$ à la correction de moyenne près
  ($-\sigma^2/2$, négligeable ici) : à $\sigma = 0{,}01$, la valeur employée
  dans tout ce programme, $\mathbb{E}|\xi| \approx 8{,}0\times10^{-3}$, donc
  un pas multiplie un capital par $1 \pm 0{,}8\,\%$ en moyenne, et par
  $1 \pm 2{,}0\,\%$ dans 5\,\% des cas. Le choc est un bruit de fond, pas un
  moteur : sur les 2000 pas d'un amorçage, sa contribution \emph{nette}
  cumulée (\code{shock\_gain}) vaut $-0{,}41\,\%$ du capital agrégé final,
  et aucun pas isolé ne déplace \code{K\_tot} de plus de $0{,}16\,\%$.
\item \textbf{Production} --- chaque entité produit selon
  \eqref{eq:production} et \emph{ajoute intégralement} sa production à son
  propre capital.
\item \textbf{Service des intérêts} --- chaque emprunteuse doit
  $\mathrm{due} = \sum_k q_k r_k$, somme sur ses contrats du principal $q_k$
  fois le taux $r_k$. Si son capital ne suffit pas, elle paie ce qu'elle
  peut, tombe à zéro et est marquée en \emph{défaut de liquidité}.
\item \textbf{Dépréciation} --- chaque capital est multiplié par $1-\delta$.
\item \textbf{Marché du crédit} --- $\lfloor \eta(N) \rfloor$ tirages de
  paires, où $N$ est le nombre d'entités éligibles et
  $\eta_{\rho,\beta}(N) = \rho\,N_{\text{ref}}(N/N_{\text{ref}})^{\beta}$
  (la baseline $\rho=1$, $\beta=1$ donne simplement $\eta(N)=N$). Dans chaque
  paire, la plus riche est la \emph{prêteuse}, la plus pauvre
  l'\emph{emprunteuse} ; le prêt ne va jamais dans l'autre sens. Un contrat
  transfère un principal $q$ et fixe un taux $r$, tous deux gelés à vie.
\item \textbf{Faillites en cascade} --- toute entité dont la valeur nette
  $K + \text{créances} - \text{dettes}$ devient négative, ou qui est en
  défaut de liquidité, meurt : ses dettes sont annulées (perte sèche pour
  ses prêteuses, qui peuvent tomber à leur tour) et son capital est détruit.
  On itère jusqu'à point fixe.
\end{enumerate}

Un identifiant d'entité n'est jamais réutilisé ; une entité morte le reste.

\subsection{Notations}

\begin{center}\footnotesize
\begin{tabular}{@{}llp{8.6cm}@{}}
\toprule
\textbf{Symbole} & \textbf{Nom} & \textbf{Définition} \\
\midrule
$K_i$ & capital & Ressource détenue par l'entité $i$. \\
$A_i$ & coefficient d'extraction & Facteur multiplicatif de la production
\eqref{eq:production}. \textbf{Levier primaire} de ce programme. \\
$\gamma_i$ & exposant d'extraction & Exposant de la production : il fixe le
rendement d'échelle. C'est lui qui rend les rendements \emph{marginaux}
décroissants ($0<\gamma<1$), et non le niveau de la production, qui croît
toujours avec le capital. Levier secondaire de ce programme. \\
$(A_i,\gamma_i)$ & technologie & Couple fixé à la naissance, jamais modifié
sauf par une intervention. Reçoit un identifiant entier interne. \\
$\lambda$ & taux de naissance & Espérance du nombre d'entités créées par pas. \\
$K_0$ & capital de naissance & Capital dont une entité neuve est dotée. \\
$\delta$ & dépréciation & Fraction du capital perdue à chaque pas. \\
$\sigma$ & amplitude du choc & Écart-type du choc multiplicatif log-normal. \\
$N$ & taille du bassin & Nombre d'entités éligibles au marché à ce pas. \\
$x,\ y$ & capitaux de la paire & $x = K_b$ (emprunteuse), $y = K_\ell$
(prêteuse), avec $y > x$ par la règle de marché. \\
$C$ & somme conservée & $C = x+y$ ; le transfert la laisse inchangée. \\
$q$ & principal & Montant transféré par le contrat. \\
$r$ & taux & Service versé par pas, perpétuel : l'emprunteuse doit $q\,r$ à
chaque pas, indéfiniment. \\
$\delta^{*}$ & transfert optimal & Principal qui maximise la production
jointe de la paire (§\ref{sec:modele}, éq.~\ref{eq:transfert}). \\
$\lambda^{*}$ & part optimale & Part de $C$ revenant à l'emprunteuse à
l'optimum ; $h(C) = C\lambda^{*}$ est son capital optimal. \\
$\Kaut$ & échelle autarcique & $[A(1-\delta)/\delta]^{1/(1-\gamma)}$ : capital
auquel production et dépréciation s'équilibrent sans marché. Repère
d'échelle du système. \\
\code{prod\_tot} & production agrégée & Somme des productions du pas.
\textbf{Variable d'intérêt de ce programme.} \\
\code{K\_tot}, \code{pop} & agrégats & Capital total et effectif vivant. \\
\code{deaths} & morts & Entités mortes au pas (racines + victimes de cascade). \\
\code{defaults} & défauts & Entités en défaut de liquidité au pas. \\
\code{roots\_insolvency} & racines d'insolvabilité & Entités dont la valeur
nette est passée sous zéro par leurs propres pertes (à distinguer des
victimes de cascade). \\
\code{loan\_volume} & volume de prêt & Somme des principaux transférés au pas. \\
\code{mkt\_rounds} & tirages & Nombre de paires tirées au pas. \\
\code{mkt\_blocked\_dir} & paires refusées & Paires où $\delta^{*} \le 0$ :
l'optimum voudrait faire circuler le capital du pauvre vers le riche, ce que
le sens du prêt interdit. \textbf{Compteur propre à M4.3Live.} \\
\bottomrule
\end{tabular}
\end{center}

\subsection{L'institution de principal de M4.3Live}

C'est le seul point où M4.3Live diffère de son prédécesseur M4.3. Là où M4.3
transférait la moitié de l'écart de capital, $q = (y-x)/2$, M4.3Live
transfère le montant qui maximise la production \emph{jointe} de la paire
immédiatement après l'échange :
\begin{equation}
\delta^{*} \;=\; \arg\max_{-x \le \delta \le y}\;
A_b\,(x+\delta)^{\gamma_b} + A_\ell\,(y-\delta)^{\gamma_\ell}
\;=\; h(C) - x,
\qquad h(C) = C\,\lambda^{*}.
\label{eq:transfert}
\end{equation}
Trois conséquences importent pour lire les deux rapports :

\begin{enumerate}
\item Quand les deux entités partagent la \emph{même} technologie,
  $\lambda^{*} = 1/2$ et $\delta^{*} = (y-x)/2$ exactement : la baseline
  homogène de M4.3Live \emph{est} celle de M4.3.
\item Dès que les technologies diffèrent, $\lambda^{*} \ne 1/2$ et le partage
  est pondéré en faveur de celle dont le \emph{rendement marginal est le plus
  élevé au capital considéré}. Il n'existe pas de « meilleure »
  technologie dans l'absolu : à exposants inégaux, l'ordre des rendements
  marginaux $\gamma A K^{\gamma-1}$ s'inverse avec le capital. Un $\gamma$
  élevé l'emporte aux grands capitaux, un $A$ élevé aux petits, et il existe
  toujours un capital de croisement. C'est pourquoi $\lambda^{*}$ dépend de
  $C$ et non seulement des paramètres.
\item $\delta^{*}$ peut être \emph{négatif}. Comme $\delta^{*} = C\lambda^{*} - x$,
  on a exactement
  \[ \delta^{*} \le 0 \iff \frac{x}{C} \ge \lambda^{*} , \]
  c'est-à-dire : dès que la prêteuse est celle dont le rendement marginal est
  le plus élevé, l'optimum voudrait lui envoyer du capital, et la paire ne
  traite pas. C'est le compteur \code{mkt\_blocked\_dir}, et c'est un canal à
  part entière, pas un détail numérique.
\end{enumerate}

\paragraph{Surplus coopératif.} Le \textbf{surplus coopératif} d'une paire
est le gain de production \emph{jointe} que l'échange produit, à capital
total inchangé :
\begin{equation}
S(x,y) \;=\;
\underbrace{A_b (x+q)^{\gamma_b} + A_\ell (y-q)^{\gamma_\ell}}
          _{\text{après le transfert}}
\;-\;
\underbrace{A_b\,x^{\gamma_b} + A_\ell\,y^{\gamma_\ell}}_{\text{avant}}
\;\ge\; 0 ,
\label{eq:surplus}
\end{equation}
où $q$ est le principal effectivement transféré --- égal à $\delta^{*}$ sauf
quand le plafond mord. La positivité est immédiate quand $q = \delta^{*}$ :
$\delta^{*}$ maximise le premier terme et $\delta = 0$ est admissible ; sous
plafond elle reste vraie car $q$ est entre $0$ et $\delta^{*}$ et le premier
terme est concave en $\delta$. $S$ est nul exactement quand la paire est
déjà à l'optimum ($x/C = \lambda^{*}$), et les paires refusées ne
contribuent pas. La colonne \code{mkt\_surplus} des séries est la somme de
$S$ sur les paires traitées au pas. C'est un \emph{flux de
production par pas}, homogène à \code{prod\_tot}, et non un stock de
capital : le marché ne crée aucun joule, il déplace du capital vers l'endroit
où il produit davantage.

\subsection{Les trois portées d'intervention}
\label{sec:portees}

Une intervention change un paramètre \emph{pendant} qu'une simulation
tourne. Sa \textbf{portée} dit qui elle touche :

\begin{description}
\item[\code{all} (toutes)] --- rétroactif \emph{et} prospectif : toutes les
  entités vivantes reçoivent la nouvelle valeur, et la valeur par défaut à la
  naissance change aussi.
\item[\code{new} (nouvelles)] --- \emph{vintage} technologique : seule la
  valeur par défaut à la naissance change. Les entités déjà vivantes gardent
  la leur pour toujours ; la population se convertit par renouvellement.
\item[\code{fraction} (une fraction $\varphi$)] --- une fraction $\varphi$
  des entités vivantes est tirée au sort une fois, et elle seule est
  modifiée. \textbf{La valeur par défaut à la naissance n'est pas touchée} :
  aucune naissance ne vient reconstituer la cohorte traitée.
\end{description}

Certaines combinaisons n'ont pas de sens et sont refusées explicitement :
$K_0$ n'existe qu'au moment de la naissance, donc \code{all} et \code{new} y
sont le même objet et \code{fraction} n'a pas d'objet ; $\lambda$, $\delta$,
$\sigma$, $\rho$ sont des paramètres de population et non des attributs
d'entité, donc \code{new} y est un alias explicite de \code{all}.


\section{Objet et périmètre}

M4.3Live succède à M4.3 (\code{m4\_3\_credit\_soc}) mais n'en est pas une
extension : l'institution de principal change de nature. Là où M4.3
transfère la moitié de l'écart de capital d'une paire, M4.3Live transfère
le montant qui \emph{maximise la production jointe} des deux entités
immédiatement après l'échange. À cela s'ajoutent deux capacités
nouvelles : des technologies $(A,\gamma)$ par entité, et la modification
de paramètres pendant qu'une simulation tourne.

Le périmètre scientifique est celui du §7 du prompt --- observer un
éventuel effet rebond --- et rien de plus. Ce rapport ne traite que la
construction ; le verdict d'observation est dans le rapport de résultats.

\section{L'institution de principal}

\subsection{Le problème}

Pour une paire (emprunteuse $b$, prêteuse $\ell$) de capitaux $x=K_b$ et
$y=K_\ell$ et de technologies $(A_b,\gamma_b)$ et $(A_\ell,\gamma_\ell)$,
le transfert $\delta$ retenu résout
\begin{equation}
\max_{-x\le\delta\le y}\; A_b\,(x+\delta)^{\gamma_b} + A_\ell\,(y-\delta)^{\gamma_\ell},
\qquad 0<\gamma<1 .
\label{eq:probleme}
\end{equation}
L'objectif est strictement concave en $\delta$ et sa dérivée diverge vers
$+\infty$ en $\delta\to-x$ comme vers $-\infty$ en $\delta\to y$ :
l'optimum est intérieur et unique. En posant $C=x+y$ (somme conservée par
l'échange) et $h(C)=C\lambda^{*}$ le capital optimal attribué à
l'emprunteuse, la solution s'écrit exactement $\delta^{*}=h(C)-x$ : à
couple de technologies fixé, toute la non-linéarité tient dans la seule
variable $C$ (rapport d'architecture §6, Hessien exactement de rang un).

\subsection{Trois régimes, routés sur des entiers}

\begin{description}
\item[(a) même technologie.] Déclenché par l'égalité des
  \emph{identifiants} \code{s\_b == s\_l} (\code{m4\_3live/kernel.py},
  \code{PrincipalKernel.solve}), et non par un test $\lambda^{*}=1/2$ sur
  des flottants. Le code applique alors la formule littérale
  $\delta=0{,}5\,(y-x)$, exactement celle de
  \code{m4\_3\_credit\_soc/m4\_3/model.py:288-289}.
\item[(b) $\gamma$ égaux, $A$ différents.] Forme fermée exacte,
  \[\lambda^{*}=\frac{A_b^{1/(1-\gamma)}}{A_b^{1/(1-\gamma)}+A_\ell^{1/(1-\gamma)}},\]
  dont le logarithme n'est payé qu'une fois, à la création du couple de
  technologies.
\item[(c) $\gamma$ différents.] Pas de forme fermée. Avec
  $u=(\gamma_\ell-\gamma_b)/(2-\gamma_b-\gamma_\ell)$ et
  $t=t_0+u\ln C$, la condition d'optimalité s'écrit
  $t=-z/2+u\,L(z)$ avec $L(z)=\ln(2\cosh(z/2))$, puis
  $\lambda^{*}=\sigma(z)$.
\end{description}

\paragraph{Le piège du régime (b), et pourquoi le routage entier compte.}
$\lambda^{*}=1/2$ \emph{si et seulement si} $A_b=A_\ell$ en plus de
$\gamma_b=\gamma_\ell$. Une intervention qui ne modifie que $A$ --- le
levier primaire de toute la campagne --- fait donc immédiatement sortir
toute paire mixte du régime (a) vers un partage \emph{pondéré}. Avec
$A_b=1{,}5$, $A_\ell=1{,}0$, $\gamma=0{,}5$ : $\lambda^{*}=0{,}692$, et
sur la paire $(x,y)=(400,1200)$ le transfert vaut $707{,}7$ contre
$400{,}0$ pour la règle arithmétique (\code{tests/test\_institution.py}).
Router le régime (a) sur $\lambda^{*}=1/2$ plutôt que sur l'identité des
technologies aurait fait de la parité avec M4.3 une coïncidence numérique
au lieu d'une propriété de construction ; le prompt (§3.6) distingue
explicitement les deux, d'où le choix retenu.

\begin{figure}[H]
\centering
\includegraphics[width=\textwidth]{figures/institution_regimes.png}
\caption{L'institution, rendue visible. À gauche, le transfert $\delta^{*}$
en fonction de la part de capital détenue par l'emprunteuse, à $C$ fixé à la
médiane du domaine réel de M4.3. La droite grise est le régime (a), qui
s'annule exactement à l'égalité des capitaux. Doper l'\emph{emprunteuse}
déplace le zéro vers la droite ; doper la \emph{prêteuse} le déplace vers la
gauche, ouvrant une bande où $\delta^{*} \le 0$ et où le marché refuse de
traiter. À droite, le rapport à la règle arithmétique de M4.3. Données :
\code{results/conception/institution\_curves.csv}.}
\label{fig:institution}
\end{figure}

\paragraph{Le seuil de refus est exactement $\lambda^{*}$.} Puisque
$\delta^{*} = C\lambda^{*} - x$, la paire ne traite pas dès que
$x/C \ge \lambda^{*}$. \fait{} Pour une prêteuse dopée $A\times1{,}5$ à
$\gamma$ égal, $\lambda^{*} = 1/(1+1{,}5^{2}) = 0{,}3077$ : toutes les paires
où l'emprunteuse détient plus de 30,8~\% du capital commun sont refusées,
soit 38,5~\% des paires si $x/C$ est uniforme sur $[0, 1/2]$. Pour une
prêteuse dopée $\gamma: 0{,}5 \to 0{,}6$, le seuil tombe à $0{,}1408$, soit
71,8~\% des paires. \inference{} C'est l'ordre de grandeur observé dans la
campagne (§7) : \code{mkt\_blocked\_dir} y atteint 22 à 37~\% des tirages
selon l'intensité du traitement.

\paragraph{Résolution du régime (c).} $|u|<1$ tant que
$\gamma\in\,]0,1[$, donc $F(z)=-z/2+uL(z)-t$ est strictement décroissante
avec $|F'|\ge(1-|u|)/2>0$ : Newton depuis $z_0=-2t$ est sûr, et
l'itération de point fixe $z\leftarrow-2t+2uL(z)$, globalement
contractante de facteur $|u|$, sert de repli dès qu'un pas de Newton
n'améliore pas $|F|$. La convergence est donc globale, sans encadrement ni
bissection. \fait{} Résidu maximal $9{,}33\cdot10^{-15}$ sur une grille
$t\in[-30,40]\times u\in[-0{,}9,0{,}9]$
(\code{tests/test\_institution.py}).

\paragraph{Vérification indépendante.} Le solveur de référence des tests
n'est pas une autre écriture de la même machinerie : c'est une bissection
sur la condition du premier ordre
$A_b\gamma_b(x+\delta)^{\gamma_b-1}-A_\ell\gamma_\ell(y-\delta)^{\gamma_\ell-1}=0$.
\fait{} Écart relatif maximal (rapporté à $C$) : $<10^{-14}$ sur 84 cas du
régime (b), $2{,}27\cdot10^{-16}$ sur 35 cas du régime (c).

\subsection{Le transfert peut être négatif --- un canal, pas un détail}

Rien ne garantit $\delta^{*}>0$. Quand la prêteuse (la plus riche, par la
règle de marché héritée, \code{m4\_3/model.py:357-361}) est celle dont le
\emph{rendement marginal} $\gamma A K^{\gamma-1}$ est le plus élevé à son
capital, l'optimum de production jointe voudrait faire circuler le capital
\emph{du pauvre vers le riche} : $\delta^{*}<0$, et le sens du prêt interdit
la transaction. Une intervention sur $A$ ne change donc pas seulement la
taille des transferts : elle \textbf{éteint le prêt} dans les paires où
l'entité dopée est la riche.

\incertitude{} On évitera de parler de « meilleure technologie » : sous
exposants inégaux il n'en existe pas dans l'absolu. L'ordre des rendements
marginaux dépend du capital --- un $\gamma$ élevé l'emporte aux grands
capitaux, un $A$ élevé aux petits --- et c'est précisément pour cela que
$\lambda^{*}$ dépend de $C$ dans le régime (c).

C'est pourquoi la série de M4.3Live sépare trois compteurs de refus là où
M4.3 n'en avait aucun : \code{mkt\_blocked\_dir} ($\delta^{*}\le0$),
\code{mkt\_blocked\_tiny} ($0<\delta^{*}<$ \code{MIN\_LOAN}) et
\code{mkt\_blocked\_rate} (taux non strictement positif, possible sous la
seule règle candidate du §\ref{sec:taux}). Sans cette séparation,
l'effondrement du volume de prêt d'un bras traité serait ininterprétable.

\subsection{Borne supérieure du transfert (décision ouverte §11)}
\label{sec:plafond}

Deux institutions étaient possibles : autoriser le transfert jusqu'à
l'optimum $h(C)$, ou le plafonner à l'égalisation des capitaux. Le choix
est \emph{économique}, pas numérique ; il ne pouvait donc pas être tranché
sur la baseline, où $\lambda^{*}=1/2$ rend les deux plafonds identiques.
Il a été piloté sur un bras \emph{traité} ($\varphi=0{,}2$, $A\times1{,}5$,
600 pas depuis $t_0=2000$, graine 0).

\fait{} Le plafond d'égalisation mord 149 fois au pas $h=1$, 24 fois à
$h=5$, une fois à $h=25$ ; au-delà de $h=50$, il ne mord plus que 23 fois en
550 pas, soit environ une transaction sur 29\,000.
\inference{} La raison est mécanique et intéressante en soi : les entités
dopées accumulent du capital et deviennent, en quelques dizaines de pas,
le côté riche de presque toutes leurs paires. Or le plafond ne peut mordre
que lorsque l'entité dopée est l'\emph{emprunteuse}. Passé ce transitoire,
les paires mixtes basculent dans le régime $\delta^{*}\le0$ décrit
ci-dessus, et \code{mkt\_blocked\_dir} se stabilise autour de 25~\% des
rounds.

\fait{} Aucun des deux plafonds ne déstabilise le système : statut
\code{ok}, $0$ défaut de liquidité dans les deux cas, morts cumulées
$18\,235$ (optimum) contre $17\,925$ (égalisation), production moyenne
$32\,672$ contre $32\,154$.

\begin{figure}[H]
\centering
\includegraphics[width=\textwidth]{figures/cap_pilot.png}
\caption{Pilote du plafond, pas à pas. À gauche, les transactions
plafonnées s'effondrent en une vingtaine de pas. À droite, la part des
tirages refusés par le sens du prêt suit la part traitée de la population :
les paires mixtes ne sont plus plafonnées, elles sont refusées. Données :
\code{results/conception/cap\_pilot\_steps.csv}.}
\label{fig:cap}
\end{figure}

\textbf{Décision : \code{transfer\_cap = "optimum"}}, l'énoncé littéral de
l'institution. Plafonner à l'égalisation reviendrait à ramener
l'institution vers la règle arithmétique précisément dans les paires où
elle en diffère --- donc à supprimer l'effet à observer --- pour un gain de
stabilité nul dans le régime mesuré. Le plafond alternatif reste
implémenté et sélectionnable en ablation.

\subsection{Le transfert ne peut pas créer d'insolvabilité dans le pas}

Le transfert baisse $K$ de $q$ et monte les créances de $q$ : la valeur
nette $K+\text{créances}-\text{dettes}$ est inchangée. Il ne peut donc pas
alimenter la cascade de faillites du même pas ; le seul risque d'une règle
non plafonnée est reporté au pas suivant, par le canal de liquidité.
\fait{} Vérifié (\code{tests/test\_institution.py}) : aucune valeur nette
négative et aucune incohérence de carnet après 40 pas.

\section{Cache du noyau : chaud, tiède, froid}
\label{sec:cache}

\subsection{Ce qui est implémenté}

Une matrice dense \code{kernel[receveuse][donneuse]} d'identifiants entiers
route un couple en $O(1)$. Les couples de régime (c) reçoivent une table
1D de $h(C)$, indexée par l'exposant \code{frexp} de $C$ (une ligne par
octave, grille uniforme en mantisse), interpolée par une cubique de
Hermite dont les pentes viennent de la formule analytique
$h'(C)=q\,h/S$, $S=p(C-h)+qh$ --- jamais d'une différence finie. Chaque
ligne est compilée paresseusement, après un seuil d'usages de son exposant.

\subsection{La table est correcte, et on le vérifie par l'ordre de convergence}

\fait{} Erreur maximale de la table contre Newton exact, mesurée sur 400
paires balayant tout le domaine réel de $C$ (quantiles 0,1~\% à 99,9~\% de
M4.3, soit 192 à 2214) : $3{,}80\cdot10^{-5}$ / $2{,}61\cdot10^{-6}$ /
$1{,}70\cdot10^{-7}$ / $1{,}05\cdot10^{-8}$ / $6{,}87\cdot10^{-10}$ pour
9 / 17 / 33 / 65 / 129 nœuds par octave, soit des rapports successifs de
$14{,}6$, $15{,}4$, $16{,}2$ et $15{,}3$. \inference{} L'ordre~4 observé
confirme que la pente analytique est juste : une dérivée fausse dégraderait
l'interpolation à l'ordre~2, et le test le détecterait
(\code{test\_lut\_convergence\_order}).

\incertitude{} Une précaution de mesure, apprise en la ratant d'abord :
l'échantillon doit \emph{balayer} $C$. Un premier jeu de 120 paires gardait
$C$ constant --- il ne sondait donc qu'un seul point de la table, dont
l'erreur dépend de sa position fortuite dans son intervalle
d'interpolation, et les rapports mesurés étaient erratiques (48,7 puis 6,3).
Les chiffres ci-dessus sont un vrai supremum sur le domaine.

Le défaut retenu est 65 nœuds, soit $1{,}05\cdot10^{-8}$ unité de capital sur
tout le domaine, pour des capitaux de l'ordre de $10^{3}$.

\begin{figure}[H]
\centering
\includegraphics[width=0.72\textwidth]{figures/kernel_convergence.png}
\caption{Erreur de la table 1D contre Newton exact, en fonction du nombre de
nœuds par octave. La pente mesurée suit l'ordre~4, pas l'ordre~2 : la dérivée
analytique $h'(C)=q\,h/S$ qui alimente l'interpolation de Hermite est donc
correcte. Le trait horizontal situe le chemin tiède. Données :
\code{results/conception/lut\_convergence.json}.}
\label{fig:convergence}
\end{figure}

\subsection{Pourquoi le chemin tiède n'est pas le défaut}

Le rapport d'architecture prescrit, pendant la phase d'observation d'un
couple, un noyau sans table : une étape de Newton depuis $z_0=-2t$
(éq.~21). Il est implémenté et testé, sous
\code{kernel\_policy="hybrid"}.

\fait{} Son erreur mesurée sur le domaine réel de M4.3 (quantiles de $C$ :
192 / 1590 / 2214) est de $9{,}8\cdot10^{-2}$ unité de capital --- du même
ordre que le $2{,}5\cdot10^{-1}$ du banc du rapport, et sept ordres de
grandeur au-dessus de la table.

\textbf{Décision : \code{kernel\_policy = "exact\_lut"} par défaut} ---
formes fermées pour (a) et (b), Newton exact pour (c) tant que la table
n'existe pas, table ensuite. Le prompt (§3.2) autorise explicitement de
s'écarter de l'architecture prescrite « sauf justification écrite d'un
meilleur choix » ; la justification est double. D'une part, aucune
approximation grossière n'entre alors jamais dans une trajectoire
scientifique : la seule erreur résiduelle est celle, contrôlée à
$1{,}5\cdot10^{-9}$, de l'interpolation. D'autre part --- et c'est
l'argument décisif --- le gain de débit du chemin tiède est sans effet
pratique, parce que le noyau ne domine pas le coût d'un pas.

\subsection{Seuil d'amortissement recalibré, et sa portée réelle}
\label{sec:seuil}

\input{tables/bench_kernel.tex}

\begin{figure}[H]
\centering
\includegraphics[width=\textwidth]{figures/kernel_cost_precision.png}
\caption{Débit et précision des chemins du noyau. Le chemin tiède est deux
fois plus lent que la table tout en étant sept ordres de grandeur moins
précis : il n'occupe aucun point utile du front. Données :
\code{results/analysis/bench\_kernel.json}.}
\label{fig:bench}
\end{figure}

\fait{} Le seuil d'amortissement mesuré sur cette machine (coût de
compilation d'une ligne divisé par l'économie par appel de la table face à
Newton exact) est de \textbf{\input{tables/bench_threshold.tex}\unskip}
usages par ligne d'exposant --- vingt-cinq fois moins que les $\sim$1800
du banc synthétique du rapport, parce que le Newton exact est ici plus
coûteux (Python scalaire, avec l'évaluation supplémentaire du garde-fou)
tandis que la compilation d'une ligne reste très bon marché
($0{,}45$\,ms).

\fait{} Mais la question ouverte du §11 se résout mieux autrement : sur un
run réel de régime (c), le noyau représente
\textbf{\input{tables/bench_share.tex}\unskip} du temps d'un pas
($1268$ appels par pas, $114$\,ms par pas). Le pas est dominé par la phase
d'intérêts, qui parcourt les $\sim$32\,000 contrats actifs du carnet.
\inference{} La valeur exacte du seuil est donc sans portée pratique :
entre 71 et 1800, l'écart de temps mural est inférieur à 1\,\% du coût
d'une campagne. Le défaut \code{lut\_threshold} est laissé à \textbf{1800},
la valeur avec laquelle toute la campagne du §7 a tourné ; l'abaisser
maintenant déplacerait le point de bascule table/Newton, donc les
trajectoires au niveau $10^{-9}$, pour un gain non mesurable. Ce qui
compte est que le chemin chaud soit \emph{exact}, pas qu'il soit rapide.

\subsection{Ce qui domine réellement le coût d'un pas}

Puisque le noyau ne pèse que 1,3~\%, il vaut la peine de dire ce qui pèse le
reste. \fait{} À population et carnet actifs constants, le coût d'un pas
\emph{croît} : $96{,}3$\,ms à $t = 210$, $102{,}9$ à $t = 610$, $111{,}5$ à
$t = 1010$, $115{,}3$ à $t = 1810$ (+20~\%), alors que le nombre de prêts
actifs reste entre 31\,000 et 37\,000 et la population autour de 1100.
\inference{} La seule quantité qui croît est le nombre de clefs du
dictionnaire \code{by\_borrower}, que la phase d'intérêts parcourt
intégralement à chaque pas : 6389 à $t = 210$, 54\,527 à $t = 1810$, dont
98~\% pointent sur un ensemble \emph{vide}. Ce sont les entités mortes :
\code{LoanBook.remove} retire l'identifiant du contrat mais ne supprime
jamais la clef (comportement hérité tel quel de
\code{m4\_3/model.py:191-202}).

\fait{} C'est mesurable et c'est assumé : purger ces clefs changerait l'ordre
d'itération de \code{by\_borrower}, donc l'ordre de sommation de
\code{interest\_paid}, donc la parité bit à bit avec M4.3 (§\ref{sec:parite}
et §\ref{sec:snapshot}, même arbitrage que pour l'ordre des ensembles).
Conséquence pratique à connaître avant de dimensionner une campagne : le coût
d'un run croît linéairement avec le nombre d'entités \emph{jamais nées},
donc avec $\lambda \times T$, et pas seulement avec la taille instantanée du
système. Données : \code{results/analysis/step\_cost\_growth.json}.

\subsection{Ce qui n'est pas implémenté, et pourquoi}

La LUT 2D globale du rapport n'est pas nécessaire ici : le support
technologique est discret, donc chaque cellule est un problème 1D
(rapport §11, même conclusion). L'« audit aléatoire » suggéré pour le
chemin froid n'est pas implémenté \emph{dans} la boucle : il consommerait
soit un tirage --- interdit, cela déplacerait la trajectoire --- soit un
générateur séparé, qui casserait le rejeu. L'audit est fait hors ligne,
dans les tests, où il compare table et Newton exact sur le domaine réel.

\section{Déterminisme}

Les trois garanties de cette section --- rejeu d'un journal, invariant de
pause, aller-retour par snapshot --- se mesurent de la même façon : l'écart
pas à pas à une trajectoire de référence. La \autoref{fig:determinisme} les
met sur un même graphique.

\begin{figure}[H]
\centering
\includegraphics[width=0.82\textwidth]{figures/determinism.png}
\caption{Les trois garanties de déterminisme, mesurées sur une même
trajectoire de 120 pas comportant deux interventions. Les trois écarts sont
\emph{exactement} nuls et reposent donc au plancher du graphique, six ordres
de grandeur sous l'epsilon machine. Données :
\code{results/conception/determinism.csv}.}
\label{fig:determinisme}
\end{figure}

\subsection{Compilation synchrone (contre la suggestion du rapport)}

Le rapport d'architecture propose de déporter la construction d'une table
sur un thread, avec remplacement atomique du pointeur (§7.2). \textbf{Cette
suggestion est refusée.} Le moment où une table remplace Newton change la
valeur retournée d'environ $10^{-9}$ ; si ce basculement dépend de
l'ordonnancement d'un thread, deux exécutions du même journal divergent, et
l'exigence de rejeu strictement identique (§4, §8) tombe. La compilation
est donc exécutée en ligne, au point de déclenchement, lui-même déterminé
par un compteur qui ne dépend que de la séquence simulée. Elle ne consomme
aucun tirage : elle change la durée murale d'un pas, jamais sa trajectoire.

\subsection{Rejeu d'un journal}

Une intervention arrivée en HTTP est mise en file, puis appliquée au tout
début du pas suivant, \emph{avant les naissances}, et journalisée avec le
$t$ auquel elle a réellement pris effet. Le rejeu soumet chaque entrée du
journal par le \emph{même} \code{Simulation.submit} : le chemin de code
vérifié est celui de la production.

\fait{} \code{tests/test\_replay.py} construit son journal à partir d'une
\emph{vraie session en direct} --- quatre interventions couvrant les trois
portées, soumises de façon asynchrone depuis le thread principal pendant
que la boucle tournait, sans savoir à quel pas elle en était. Rejoué sans
tête, ce journal reproduit les 120 pas à l'identique sur 16 colonnes, avec
les mêmes $t$ effectifs \emph{et} les mêmes identifiants tirés.

\subsection{Pause et pas-à-pas}

La pause est une porte placée \emph{avant} \code{step()} : elle ne touche
pas au générateur. \fait{} 80 pas entrecoupés de quatre pauses franches, et
30 pas déclenchés un par un, produisent des trajectoires bit-identiques à
une lecture continue.

\subsection{Snapshots : une réserve mesurée et bornée}
\label{sec:snapshot}

Un snapshot emporte la population, le carnet, l'état du générateur, le
registre de technologies \emph{et les compteurs d'usage du noyau} --- sans
ces derniers, une branche restaurée compilerait ses tables à un autre
indice d'appel et divergerait de la branche continuée.

\fait{} Un aller-retour laisse la \textbf{dynamique bit-identique}
($K_{\mathrm{tot}}$ à $0$ près sur 90 pas), mais \code{interest\_paid}
peut différer au dernier bit ($2{,}1\cdot10^{-16}$ relatif).
\inference{} Cause identifiée précisément : la phase d'intérêts parcourt
\code{book.by\_borrower[borrower]}, un \code{set}, et CPython ne sérialise
pas la disposition interne d'un ensemble ; après \code{pickle}, l'ordre
d'itération change pour 225 des 467 ensembles mesurés. Cet ordre n'affecte
\emph{aucun} capital --- chaque versement va à une prêteuse différente,
donc dans un accumulateur différent --- seulement deux sommes de
diagnostic, \code{interest\_paid} et \code{int\_out}.

\begin{figure}[H]
\centering
\includegraphics[width=0.72\textwidth]{figures/snapshot_round_trip.png}
\caption{Aller-retour par snapshot : écart à une simulation menée d'un
trait. Le capital et la production restent \emph{exactement} nuls (tracés au
plancher du graphique) ; seul \code{interest\_paid} bouge, au niveau de
l'epsilon machine. Données :
\code{results/conception/snapshot\_round\_trip.csv}.}
\label{fig:snapshot}
\end{figure}

La réserve est sans conséquence pour le protocole : les bras d'une même
graine sont chargés depuis le \emph{même} fichier, donc partagent la même
disposition d'ensembles et coïncident exactement jusqu'à leur
intervention. La corriger exigerait de remplacer les \code{set} par des
structures ordonnées, donc de changer l'ordre de sommation hérité de M4.3
et de perdre la parité. Elle est documentée plutôt que corrigée.

\section{Fork indépendant et parité}

\subsection{Indépendance}

Aucun module de \code{m4\_3live\_credit\_soc} n'importe un autre moteur du
dépôt. \code{m4\_3\_credit\_soc/m4\_3/model.py} est une \emph{référence
structurelle lue}, pas une dépendance : listes parallèles jamais
réindexées, carnet à agrégats incrémentaux, RNG
\code{numpy.random.default\_rng(seed)}, ordre des phases, faillites
cancel+destroy, pool $k\equiv2$ et $\eta_{\rho,\beta}(N)$ ont été
reproduits à la lecture.

\subsection{Parité, obtenue et vérifiée}
\label{sec:parite}

\fait{} Un run homogène de M4.3Live (graine 0, paramètres baseline,
$T=8000$) reproduit \textbf{bit à bit} le run stocké
\code{m4\_3\_credit\_soc} \code{m4\_3\_\_d1\_\_baseline\_\_seed0} :
26 colonnes de série, 8000 pas, écart maximal nul --- pas même au dernier
bit. Le noyau a servi $9\,366\,225$ fois, toutes sur le chemin identité.
\inference{} Contrôle croisé : le rapport d'architecture dénombre
$9\,366\,177$ événements de prêt sur ce même run, soit exactement 48 de
moins ; ce sont les paires dont le principal tombe sous \code{MIN\_LOAN} et
qui ne donnent donc pas de contrat.

\begin{figure}[H]
\centering
\includegraphics[width=\textwidth]{figures/parity_m4_3.png}
\caption{Parité bit à bit. En haut, la trajectoire de M4.3Live en régime
homogène superposée à celle du run M4.3 stocké --- les deux courbes sont la
même. En bas, l'écart maximal pas à pas sur les 26 colonnes : il vaut
exactement zéro, donc la courbe repose sur le plancher du graphique, deux
ordres de grandeur sous l'epsilon machine. Données :
\code{results/analysis/parity\_deviations\_8000.csv}.}
\label{fig:parite}
\end{figure}

Le test est dans \code{tests/test\_parity\_m4\_3.py} (500 pas par défaut,
les 8000 pas avec \code{--full}) : la mesure citée ici est reproductible
depuis le dépôt, elle ne vit pas dans un script jetable.

Le prompt (§3.6) n'exigeait plus cette parité. Elle a été conservée parce
qu'elle est gratuite --- elle découle du routage entier --- et parce
qu'elle vérifie d'un coup l'ordre des phases, la séquence exacte d'appels
au générateur (Poisson, normale, puis le chemin rapide $k=2$ de
l'échantillonnage) et la fonction de taux.

\section{Sémantique des trois portées}

\begin{description}
\item[\code{all}] --- rétroactif \emph{et} prospectif : toutes les entités
  vivantes reçoivent la nouvelle valeur, et la valeur par défaut à la
  naissance est modifiée.
\item[\code{new}] --- vintage technologique : seule la valeur par défaut à
  la naissance change ; les entités déjà vivantes gardent la leur pour
  toujours.
\item[\code{fraction}] --- tirage unique de $\varphi$ des vivantes, avec le
  générateur de la simulation. \textbf{La valeur par défaut à la naissance
  n'est pas modifiée} : une entité née après $t_0$ reçoit l'ancienne
  valeur. C'est ce qui garde l'intensité de traitement fixée au moment de
  la mesure ; si \code{fraction} contaminait le défaut, elle se comporterait
  comme \code{new} plus un tirage ponctuel. Vérifié explicitement
  (\code{tests/test\_scopes.py}).
\end{description}

Les technologies se composent par coordonnée : une intervention \code{all}
sur $A$ après une \code{fraction} sur $\gamma$ laisse les $\gamma$
hétérogènes intacts, chaque entité recevant $(A_{\text{nouveau}},
\gamma_i)$.

\paragraph{Dégénérescences, affichées et non masquées.} \code{K0} n'existe
qu'à la naissance : \code{all} et \code{new} sont identiques par
construction (vérifié : trajectoires bit-identiques), et \code{fraction}
n'a pas de sens. \code{lam}, \code{delta}, \code{sigma}, \code{rho},
\code{eta\_beta}, \code{eta\_n\_ref} sont des paramètres de
population/marché : \code{new} est un alias explicite de \code{all}
(vérifié), \code{fraction} n'a pas de sens. Ces portées sans objet sont
\emph{refusées} avec un message, et l'IHM affiche le texte de dégénérescence
sous le sélecteur au lieu de le désactiver en silence. \code{target\_rule}
est exclu de l'intervention : M4.3Live n'a qu'une institution de principal.

\paragraph{Co-variation de $K_0$ (décision ouverte §11).} $K_0$ n'est
\textbf{pas} co-varié automatiquement avec une intervention sur
$A$/$\gamma$. Raison : le protocole principal agit par la portée
\code{fraction} sur des entités \emph{déjà vivantes}, dont le capital est
déterminé par leur histoire ; $K_0$ ne concerne que les naissances futures,
identiques dans le bras traité et dans sa référence appariée. Co-varier
introduirait un second changement simultané --- exactement la confusion que
le prompt met en garde de créer. \incertitude{} Le caveat demeure et est
signalé dans le rapport de résultats : une intervention sur $A$ déplace
l'échelle autarcique $K^{*}_{\mathrm{aut}}=[A(1-\delta)/\delta]^{1/(1-\gamma)}
\propto A^{1/(1-\gamma)}$, donc change la tension du système en même temps
que l'extraction. L'option de co-variation reste disponible en émettant une
intervention \code{K0} distincte.

\section{Architecture « en direct »}

\subsection{Un thread, pas un sous-processus}

\code{simulation\_lab/jobs.py} lance ses runs batch en
\code{multiprocessing.Process} (\code{jobs.py:24-51}). Ce modèle est
inadapté ici : l'IHM interroge l'état toutes les demi-secondes et la
population peut atteindre \code{pop\_max}$=30\,000$ entités ; faire
transiter cet état par une \code{Queue} à chaque pas coûterait plus que le
pas. Un \code{threading.Thread} partage l'état sans copie.

\subsection{Aucun verrou en lecture}

Le thread de boucle est le seul écrivain, et il n'\emph{ajoute} que des
dictionnaires déjà complets à \code{series} / \code{tech\_series}.
\code{list.append} et le découpage de liste étant atomiques sous le GIL, un
lecteur ne peut jamais voir un enregistrement partiel. Seule la file
d'interventions, écrite par un thread HTTP et lue par le thread de boucle,
est protégée par un verrou.

\subsection{Un point d'entrée unique}

\code{Simulation.submit} est appelé aussi bien par l'IHM que par le pilote
sans tête. C'est ce qui rend le test de rejeu probant : il ne compare pas
un plan à lui-même, il rejoue le journal d'une session réelle.

\subsection{Mécanisme de sondage (décision ouverte §11)}

La convention existante de \code{simulation\_lab} (sondage HTTP périodique,
cf. \code{web/static/app.js}) est reprise, avec deux adaptations dictées
par le débit d'un direct : cadence de 500\,ms au lieu du rythme des jobs,
et transfert \emph{incrémental} --- le client envoie \code{?since=<t>} et
ne reçoit que les lignes postérieures, qu'il accumule localement. La charge
utile reste ainsi bornée quelle que soit la longueur du run. Aucune
dépendance nouvelle : \code{http.server}, JSON, et des graphiques dessinés
en \code{canvas} 2D.

\section{Intégration à simulation\_lab}

\subsection{Isolation (décision ouverte §11)}

Le choix était entre étendre le serveur existant et en lancer un second.
\textbf{Décision : un seul serveur}, donc une seule adresse pour
l'utilisateur, et un accès direct en lecture seule à
\code{RunStorage} pour la reprise. L'isolation est obtenue autrement que
par la séparation des processus : tout le code de M4.3Live vit dans son
propre dossier (\code{m4\_3live\_credit\_soc/web/router.py}), et
\code{simulation\_lab} ne gagne qu'un module minuscule
(\code{simulation\_lab/live/}) plus deux lignes en tête de \code{do\_GET} et
\code{do\_POST} (\code{web/app.py}). Aucune branche existante n'est
modifiée. Le routeur est chargé \emph{par chemin de fichier}, ce qui évite
d'introduire un nom de premier niveau générique (\code{web}) dans
\code{sys.modules}, et son chargement est paresseux : si M4.3Live est
absent ou cassé, \code{/live} renvoie 503 et le reste du laboratoire
fonctionne exactement comme avant.

\subsection{Non-régression}

\input{tables/nonregression.tex}

\section{Reprise depuis un run M4.3 et vérification de divergence}

Une session peut repartir d'un run M4.3 déjà stocké : ses paramètres et sa
graine sont lus dans les métadonnées, et le rejeu jusqu'à $t_0$ emprunte la
boucle \emph{normale}, ce qui rend la progression visible dans l'IHM.

M4.3Live étant un fork d'institution différente, rejouer un run stocké ne
reproduit sa trajectoire que dans le régime homogène --- et rien ne le
garantit hors de ce régime. Une reprise n'est donc \textbf{jamais}
présentée comme identique sans mesure : dès que $t_0$ est atteint, le
sondage d'état calcule et renvoie un rapport de divergence (nombre de pas
comparés, identité bit à bit, premier pas fautif, écarts absolus par
colonne, écart relatif maximal), écrit aussi dans
\code{divergence.json}.

\begin{figure}[H]
\centering
\includegraphics[width=0.72\textwidth]{figures/divergence_cases.png}
\caption{Le rapport de divergence mis à l'épreuve sur trois cas. Il rend
zéro quand la reprise est exacte, et il dénonce l'écart --- avec le pas et la
colonne fautifs --- quand elle ne l'est pas. Données :
\code{results/conception/divergence\_cases.json}.}
\label{fig:divergence}
\end{figure}

\fait{} Résultats (\code{tests/test\_resume\_divergence.py}) :
\begin{itemize}
\item reprise homogène de \code{m4\_3\_\_d1\_\_baseline\_\_seed0} à
  $t_0=400$ : écart maximal \emph{nul} sur 400 pas $\times$ 10 colonnes ;
\item falsification --- confronté à la série d'un \emph{autre} run
  (\code{rho\_2}), le même rapport signale la divergence dès $t=1$ sur
  \code{n\_loans}, écart relatif maximal $5{,}79\cdot10^{-1}$. Le rapport
  n'est donc pas un « toujours identique » décoratif ;
\item après une intervention à $t=61$, la divergence est détectée à partir
  de $t=61$ exactement, écart relatif maximal $10{,}0$ ;
\item le pilote sans tête imprime le rapport \emph{et} l'écrit dans
  \code{summary.json} --- il est rapporté, pas seulement calculé.
\end{itemize}

\section{Chantier taux / partage du surplus}
\label{sec:taux}

Le prompt (§3.4) ouvre une piste : faire que les entités se partagent le
surplus coopératif au lieu d'utiliser le rendement marginal, via un
paramètre de partage $p$. Il identifie correctement le vrai obstacle : ce
n'est pas le partage, c'est la conversion d'une quantité en un taux
perpétuel.

\paragraph{L'hypothèse manquante, et notre lecture.} Le prompt décrit
$\Delta$ comme « une quantité ponctuelle ». \hyp{} Sur les termes du
modèle, ce n'en est pas une : la production $A K^{\gamma}$ est un
\emph{flux par pas}, donc le surplus qu'engendre la réallocation est lui
aussi un flux par pas, qui persiste tant que l'allocation persiste. On le
gèle au moment du contrat --- exactement comme le taux marginal
$m=A\gamma K^{\gamma-1}$ est déjà gelé au contrat dans M4.3 --- et on
impose que le service perpétuel verse à la prêteuse sa part du surplus par
pas :
\[ r\,q = p\,\Delta \qquad\Longrightarrow\qquad r = \frac{p\,\Delta}{q}. \]
C'est cette convention (« rente perpétuelle au surplus gelé ») qui rend le
problème bien posé, sans horizon ni actualisation implicites. Si l'on
tenait à la lecture ponctuelle du prompt, il faudrait au contraire
introduire un taux d'actualisation arbitraire ; on a préféré expliciter la
lecture qui n'en a pas besoin.

\begin{figure}[H]
\centering
\includegraphics[width=\textwidth]{figures/rate_rules.png}
\caption{Les deux règles de taux sur le domaine réel. La règle candidate
$r = p\Delta/q$ reste partout sous le taux marginal : elle allège le service
au lieu de l'alourdir. Données : \code{results/conception/rate\_rules.csv}.}
\label{fig:taux}
\end{figure}

\paragraph{Vérifications exigées.} \fait{} $r>0$ sur les 17 paires
représentatives testées, et $\Delta\le0$ interdit le contrat (un taux nul
serait rejeté par le contrôle de cohérence du carnet). \fait{} Le rapport
$r/r_{\text{marginal}}$ vaut au plus $0{,}414$ sur un balayage complet du
domaine (\autoref{fig:taux}) et $0{,}227$ sur les paires du test : le service est plus
\emph{léger}, pas plus lourd --- le risque de pic de défauts mécanique
signalé par le prompt ne se matérialise pas. \fait{} Sur 120 pas à graine
identique, $0$ défaut sous les deux règles, intérêts versés
$3\,011\,813$ (marginal) contre $295\,574$ (surplus).

\textbf{Décision : chantier abouti mais désactivé par défaut.}
\code{rate\_rule="marginal"} reste le mécanisme opérationnel, conformément
au §3.4, et aucune cellule de la campagne n'utilise la règle candidate : la
question du rebond est mesurable avec le mécanisme existant, et mélanger
deux changements institutionnels rendrait le résultat illisible.

\section{Récapitulatif des décisions laissées ouvertes (§11)}

\begin{longtable}{@{}p{4.6cm}p{5.2cm}p{6.2cm}@{}}
\toprule
\textbf{Point ouvert} & \textbf{Décision} & \textbf{Motif} \\
\midrule
\endhead
Nom du paquet & \code{m4\_3live\_credit\_soc} / \code{m4\_3live}, inchangé
& Aucune gêne rencontrée. \\
Serveur dédié vs extension & Extension isolée du serveur existant
& Une seule adresse ; \code{RunStorage} accessible pour la reprise ;
isolation obtenue par le dossier, pas par le processus (§7.1). \\
Mécanisme de sondage & Convention \code{app.js}, cadence 500\,ms,
transfert incrémental \code{?since=} & Charge utile bornée quelle que soit
la longueur du run. \\
Hétérogénéité au-delà de $A/\gamma$ & Non & Hors mandat par défaut ; le
protocole §7 n'en a pas révélé le besoin. \\
Co-variation de $K_0$ & Délibérément découplée, option manuelle
& La portée \code{fraction} agit sur des entités dont $K$ est déjà
déterminé ; co-varier ajouterait un second changement simultané. \\
Borne du transfert & \code{optimum} (défaut), \code{equalization} en
ablation & Piloté sur un bras traité : le plafond ne mord que $h\le25$,
jamais ensuite (§\ref{sec:plafond}). \\
Seuils chaud/tiède/froid & \code{exact\_lut} par défaut ; seuil mesuré à
71, laissé à 1800 & L'erreur du chemin tiède est $10^{7}\times$ celle de la
table, pour un gain sans effet : le noyau ne pèse que 1,3\,\% du coût d'un
pas (§\ref{sec:seuil}). \\
Chantier taux/$p$ & Abouti, hypothèse documentée, désactivé par défaut
& Ne pas mélanger deux changements institutionnels dans la mesure du
rebond (§\ref{sec:taux}). \\
\bottomrule
\end{longtable}

\section{Tests}

Assertions Python simples, convention du dépôt (pas de pytest).

\begin{longtable}{@{}p{4.5cm}p{11.5cm}@{}}
\toprule
\textbf{Fichier} & \textbf{Ce qui est vérifié} \\
\midrule
\endhead
\code{test\_institution.py} & Les trois régimes contre une bissection
indépendante ; $\lambda^{*}\ne1/2$ dès que $A_b\ne A_\ell$ ; ordre de
convergence 4 de la table ; erreur réelle du chemin tiède ; $\delta^{*}$
est bien un maximum ; valeur nette préservée ; le compteur de refus par le
sens du prêt se déclenche. \\
\code{test\_scopes.py} & Les trois portées, la non-contamination du défaut
de naissance par \code{fraction}, le déterminisme du tirage, les ids
explicites sans tirage, les dégénérescences \code{K0} et population,
le refus des portées sans objet, l'ordre intervention/naissances, la
composition de deux interventions. \\
\code{test\_replay.py} & Rejeu du journal d'une \emph{vraie} session en
direct ; invariant de pause ; pas-à-pas ; aller-retour snapshot ;
restauration du cache du noyau. \\
\code{test\_resume\_divergence.py} & Reprise homogène bit à bit ;
falsification du rapport ; divergence après intervention ; remontée du
rapport jusqu'à l'appelant. \\
\code{test\_surplus\_rate.py} & $r>0$, $r<r_{\text{marginal}}$, linéarité
en $p$, blocage si $\Delta\le0$, absence de pic de défauts. \\
\code{test\_parity\_m4\_3.py} & Parité bit à bit avec le run M4.3 stocké,
sur les 26 colonnes de série, sans aucune tolérance ; 500 pas par défaut,
les 8000 pas cités dans ce rapport avec \code{--full}. \\
\bottomrule
\end{longtable}

\appendix
\section{Annexe de traçabilité : les simulations citées}

Toute simulation mentionnée dans ce rapport est identifiable par son
\code{run\_id} et consultable dans \code{simulation\_lab}. Le tableau
ci-dessous couvre les deux rapports du programme ; les figures et tableaux
référencés dans la colonne « apparaît dans » sont ceux du rapport de
résultats sauf mention contraire.

{\scriptsize
\setlength{\tabcolsep}{4pt}
\begin{longtable}{@{}>{\raggedright\arraybackslash}p{5.3cm}>{\raggedright\arraybackslash}p{1.9cm}>{\raggedright\arraybackslash}p{3.2cm}>{\raggedright\arraybackslash}p{4.9cm}@{}}
\toprule
\textbf{run\_id (simulation\_lab)} & \textbf{groupe} & \textbf{apparaît dans} & \textbf{rôle / interventions} \\
\midrule
\endfirsthead
\toprule
\textbf{run\_id (simulation\_lab)} & \textbf{groupe} & \textbf{apparaît dans} & \textbf{rôle / interventions} \\
\midrule
\endhead
\code{m4\_3live\_v2\_\_\allowbreak{}arm\_\_\allowbreak{}free\_\_\allowbreak{}all\_A150\_\_\allowbreak{}seed0} & lot D & campagne A/B & hausse globale $A\textbackslash{}times1{,}5$ : reste homogène, donc insensible à la règle de sens \\ \textit{t=2001 A$\rightarrow$1.5 (all)} \\
\code{m4\_3live\_v2\_\_\allowbreak{}arm\_\_\allowbreak{}free\_\_\allowbreak{}all\_A150\_\_\allowbreak{}seed1} & lot D & campagne A/B & hausse globale $A\textbackslash{}times1{,}5$ : reste homogène, donc insensible à la règle de sens \\ \textit{t=2001 A$\rightarrow$1.5 (all)} \\
\code{m4\_3live\_v2\_\_\allowbreak{}arm\_\_\allowbreak{}free\_\_\allowbreak{}all\_A150\_\_\allowbreak{}seed10} & lot D & campagne A/B & hausse globale $A\textbackslash{}times1{,}5$ : reste homogène, donc insensible à la règle de sens \\ \textit{t=2001 A$\rightarrow$1.5 (all)} \\
\code{m4\_3live\_v2\_\_\allowbreak{}arm\_\_\allowbreak{}free\_\_\allowbreak{}all\_A150\_\_\allowbreak{}seed11} & lot D & campagne A/B & hausse globale $A\textbackslash{}times1{,}5$ : reste homogène, donc insensible à la règle de sens \\ \textit{t=2001 A$\rightarrow$1.5 (all)} \\
\code{m4\_3live\_v2\_\_\allowbreak{}arm\_\_\allowbreak{}free\_\_\allowbreak{}all\_A150\_\_\allowbreak{}seed2} & lot D & campagne A/B & hausse globale $A\textbackslash{}times1{,}5$ : reste homogène, donc insensible à la règle de sens \\ \textit{t=2001 A$\rightarrow$1.5 (all)} \\
\code{m4\_3live\_v2\_\_\allowbreak{}arm\_\_\allowbreak{}free\_\_\allowbreak{}all\_A150\_\_\allowbreak{}seed3} & lot D & campagne A/B & hausse globale $A\textbackslash{}times1{,}5$ : reste homogène, donc insensible à la règle de sens \\ \textit{t=2001 A$\rightarrow$1.5 (all)} \\
\code{m4\_3live\_v2\_\_\allowbreak{}arm\_\_\allowbreak{}free\_\_\allowbreak{}all\_A150\_\_\allowbreak{}seed4} & lot D & campagne A/B & hausse globale $A\textbackslash{}times1{,}5$ : reste homogène, donc insensible à la règle de sens \\ \textit{t=2001 A$\rightarrow$1.5 (all)} \\
\code{m4\_3live\_v2\_\_\allowbreak{}arm\_\_\allowbreak{}free\_\_\allowbreak{}all\_A150\_\_\allowbreak{}seed5} & lot D & campagne A/B & hausse globale $A\textbackslash{}times1{,}5$ : reste homogène, donc insensible à la règle de sens \\ \textit{t=2001 A$\rightarrow$1.5 (all)} \\
\code{m4\_3live\_v2\_\_\allowbreak{}arm\_\_\allowbreak{}free\_\_\allowbreak{}all\_A150\_\_\allowbreak{}seed6} & lot D & campagne A/B & hausse globale $A\textbackslash{}times1{,}5$ : reste homogène, donc insensible à la règle de sens \\ \textit{t=2001 A$\rightarrow$1.5 (all)} \\
\code{m4\_3live\_v2\_\_\allowbreak{}arm\_\_\allowbreak{}free\_\_\allowbreak{}all\_A150\_\_\allowbreak{}seed7} & lot D & campagne A/B & hausse globale $A\textbackslash{}times1{,}5$ : reste homogène, donc insensible à la règle de sens \\ \textit{t=2001 A$\rightarrow$1.5 (all)} \\
\code{m4\_3live\_v2\_\_\allowbreak{}arm\_\_\allowbreak{}free\_\_\allowbreak{}all\_A150\_\_\allowbreak{}seed8} & lot D & campagne A/B & hausse globale $A\textbackslash{}times1{,}5$ : reste homogène, donc insensible à la règle de sens \\ \textit{t=2001 A$\rightarrow$1.5 (all)} \\
\code{m4\_3live\_v2\_\_\allowbreak{}arm\_\_\allowbreak{}free\_\_\allowbreak{}all\_A150\_\_\allowbreak{}seed9} & lot D & campagne A/B & hausse globale $A\textbackslash{}times1{,}5$ : reste homogène, donc insensible à la règle de sens \\ \textit{t=2001 A$\rightarrow$1.5 (all)} \\
\code{m4\_3live\_v2\_\_\allowbreak{}arm\_\_\allowbreak{}free\_\_\allowbreak{}control\_\_\allowbreak{}seed0} & lot D & campagne A/B ; contrôle de stationnarité & référence inerte sous sens libre \\
\code{m4\_3live\_v2\_\_\allowbreak{}arm\_\_\allowbreak{}free\_\_\allowbreak{}control\_\_\allowbreak{}seed1} & lot D & campagne A/B ; contrôle de stationnarité & référence inerte sous sens libre \\
\code{m4\_3live\_v2\_\_\allowbreak{}arm\_\_\allowbreak{}free\_\_\allowbreak{}control\_\_\allowbreak{}seed10} & lot D & campagne A/B ; contrôle de stationnarité & référence inerte sous sens libre \\
\code{m4\_3live\_v2\_\_\allowbreak{}arm\_\_\allowbreak{}free\_\_\allowbreak{}control\_\_\allowbreak{}seed11} & lot D & campagne A/B ; contrôle de stationnarité & référence inerte sous sens libre \\
\code{m4\_3live\_v2\_\_\allowbreak{}arm\_\_\allowbreak{}free\_\_\allowbreak{}control\_\_\allowbreak{}seed2} & lot D & campagne A/B ; contrôle de stationnarité & référence inerte sous sens libre \\
\code{m4\_3live\_v2\_\_\allowbreak{}arm\_\_\allowbreak{}free\_\_\allowbreak{}control\_\_\allowbreak{}seed3} & lot D & campagne A/B ; contrôle de stationnarité & référence inerte sous sens libre \\
\code{m4\_3live\_v2\_\_\allowbreak{}arm\_\_\allowbreak{}free\_\_\allowbreak{}control\_\_\allowbreak{}seed4} & lot D & campagne A/B ; contrôle de stationnarité & référence inerte sous sens libre \\
\code{m4\_3live\_v2\_\_\allowbreak{}arm\_\_\allowbreak{}free\_\_\allowbreak{}control\_\_\allowbreak{}seed5} & lot D & campagne A/B ; contrôle de stationnarité & référence inerte sous sens libre \\
\code{m4\_3live\_v2\_\_\allowbreak{}arm\_\_\allowbreak{}free\_\_\allowbreak{}control\_\_\allowbreak{}seed6} & lot D & campagne A/B ; contrôle de stationnarité & référence inerte sous sens libre \\
\code{m4\_3live\_v2\_\_\allowbreak{}arm\_\_\allowbreak{}free\_\_\allowbreak{}control\_\_\allowbreak{}seed7} & lot D & campagne A/B ; contrôle de stationnarité & référence inerte sous sens libre \\
\code{m4\_3live\_v2\_\_\allowbreak{}arm\_\_\allowbreak{}free\_\_\allowbreak{}control\_\_\allowbreak{}seed8} & lot D & campagne A/B ; contrôle de stationnarité & référence inerte sous sens libre \\
\code{m4\_3live\_v2\_\_\allowbreak{}arm\_\_\allowbreak{}free\_\_\allowbreak{}control\_\_\allowbreak{}seed9} & lot D & campagne A/B ; contrôle de stationnarité & référence inerte sous sens libre \\
\code{m4\_3live\_v2\_\_\allowbreak{}arm\_\_\allowbreak{}free\_\_\allowbreak{}new\_A075\_\_\allowbreak{}seed0} & lot D & campagne A/B & vintage $A\textbackslash{}times0{,}75$, sens libre \\ \textit{t=2001 A$\rightarrow$0.75 (new)} \\
\code{m4\_3live\_v2\_\_\allowbreak{}arm\_\_\allowbreak{}free\_\_\allowbreak{}new\_A075\_\_\allowbreak{}seed1} & lot D & campagne A/B & vintage $A\textbackslash{}times0{,}75$, sens libre \\ \textit{t=2001 A$\rightarrow$0.75 (new)} \\
\code{m4\_3live\_v2\_\_\allowbreak{}arm\_\_\allowbreak{}free\_\_\allowbreak{}new\_A075\_\_\allowbreak{}seed10} & lot D & campagne A/B & vintage $A\textbackslash{}times0{,}75$, sens libre \\ \textit{t=2001 A$\rightarrow$0.75 (new)} \\
\code{m4\_3live\_v2\_\_\allowbreak{}arm\_\_\allowbreak{}free\_\_\allowbreak{}new\_A075\_\_\allowbreak{}seed11} & lot D & campagne A/B & vintage $A\textbackslash{}times0{,}75$, sens libre \\ \textit{t=2001 A$\rightarrow$0.75 (new)} \\
\code{m4\_3live\_v2\_\_\allowbreak{}arm\_\_\allowbreak{}free\_\_\allowbreak{}new\_A075\_\_\allowbreak{}seed2} & lot D & campagne A/B & vintage $A\textbackslash{}times0{,}75$, sens libre \\ \textit{t=2001 A$\rightarrow$0.75 (new)} \\
\code{m4\_3live\_v2\_\_\allowbreak{}arm\_\_\allowbreak{}free\_\_\allowbreak{}new\_A075\_\_\allowbreak{}seed3} & lot D & campagne A/B & vintage $A\textbackslash{}times0{,}75$, sens libre \\ \textit{t=2001 A$\rightarrow$0.75 (new)} \\
\code{m4\_3live\_v2\_\_\allowbreak{}arm\_\_\allowbreak{}free\_\_\allowbreak{}new\_A075\_\_\allowbreak{}seed4} & lot D & campagne A/B & vintage $A\textbackslash{}times0{,}75$, sens libre \\ \textit{t=2001 A$\rightarrow$0.75 (new)} \\
\code{m4\_3live\_v2\_\_\allowbreak{}arm\_\_\allowbreak{}free\_\_\allowbreak{}new\_A075\_\_\allowbreak{}seed5} & lot D & campagne A/B & vintage $A\textbackslash{}times0{,}75$, sens libre \\ \textit{t=2001 A$\rightarrow$0.75 (new)} \\
\code{m4\_3live\_v2\_\_\allowbreak{}arm\_\_\allowbreak{}free\_\_\allowbreak{}new\_A075\_\_\allowbreak{}seed6} & lot D & campagne A/B & vintage $A\textbackslash{}times0{,}75$, sens libre \\ \textit{t=2001 A$\rightarrow$0.75 (new)} \\
\code{m4\_3live\_v2\_\_\allowbreak{}arm\_\_\allowbreak{}free\_\_\allowbreak{}new\_A075\_\_\allowbreak{}seed7} & lot D & campagne A/B & vintage $A\textbackslash{}times0{,}75$, sens libre \\ \textit{t=2001 A$\rightarrow$0.75 (new)} \\
\code{m4\_3live\_v2\_\_\allowbreak{}arm\_\_\allowbreak{}free\_\_\allowbreak{}new\_A075\_\_\allowbreak{}seed8} & lot D & campagne A/B & vintage $A\textbackslash{}times0{,}75$, sens libre \\ \textit{t=2001 A$\rightarrow$0.75 (new)} \\
\code{m4\_3live\_v2\_\_\allowbreak{}arm\_\_\allowbreak{}free\_\_\allowbreak{}new\_A075\_\_\allowbreak{}seed9} & lot D & campagne A/B & vintage $A\textbackslash{}times0{,}75$, sens libre \\ \textit{t=2001 A$\rightarrow$0.75 (new)} \\
\code{m4\_3live\_v2\_\_\allowbreak{}arm\_\_\allowbreak{}free\_\_\allowbreak{}new\_A150\_\_\allowbreak{}seed0} & lot D & campagne A/B ; régime nouveau & vintage $A\textbackslash{}times1{,}5$, sens libre \\ \textit{t=2001 A$\rightarrow$1.5 (new)} \\
\code{m4\_3live\_v2\_\_\allowbreak{}arm\_\_\allowbreak{}free\_\_\allowbreak{}new\_A150\_\_\allowbreak{}seed1} & lot D & campagne A/B ; régime nouveau & vintage $A\textbackslash{}times1{,}5$, sens libre \\ \textit{t=2001 A$\rightarrow$1.5 (new)} \\
\code{m4\_3live\_v2\_\_\allowbreak{}arm\_\_\allowbreak{}free\_\_\allowbreak{}new\_A150\_\_\allowbreak{}seed10} & lot D & campagne A/B ; régime nouveau & vintage $A\textbackslash{}times1{,}5$, sens libre \\ \textit{t=2001 A$\rightarrow$1.5 (new)} \\
\code{m4\_3live\_v2\_\_\allowbreak{}arm\_\_\allowbreak{}free\_\_\allowbreak{}new\_A150\_\_\allowbreak{}seed11} & lot D & campagne A/B ; régime nouveau & vintage $A\textbackslash{}times1{,}5$, sens libre \\ \textit{t=2001 A$\rightarrow$1.5 (new)} \\
\code{m4\_3live\_v2\_\_\allowbreak{}arm\_\_\allowbreak{}free\_\_\allowbreak{}new\_A150\_\_\allowbreak{}seed2} & lot D & campagne A/B ; régime nouveau & vintage $A\textbackslash{}times1{,}5$, sens libre \\ \textit{t=2001 A$\rightarrow$1.5 (new)} \\
\code{m4\_3live\_v2\_\_\allowbreak{}arm\_\_\allowbreak{}free\_\_\allowbreak{}new\_A150\_\_\allowbreak{}seed3} & lot D & campagne A/B ; régime nouveau & vintage $A\textbackslash{}times1{,}5$, sens libre \\ \textit{t=2001 A$\rightarrow$1.5 (new)} \\
\code{m4\_3live\_v2\_\_\allowbreak{}arm\_\_\allowbreak{}free\_\_\allowbreak{}new\_A150\_\_\allowbreak{}seed4} & lot D & campagne A/B ; régime nouveau & vintage $A\textbackslash{}times1{,}5$, sens libre \\ \textit{t=2001 A$\rightarrow$1.5 (new)} \\
\code{m4\_3live\_v2\_\_\allowbreak{}arm\_\_\allowbreak{}free\_\_\allowbreak{}new\_A150\_\_\allowbreak{}seed5} & lot D & campagne A/B ; régime nouveau & vintage $A\textbackslash{}times1{,}5$, sens libre \\ \textit{t=2001 A$\rightarrow$1.5 (new)} \\
\code{m4\_3live\_v2\_\_\allowbreak{}arm\_\_\allowbreak{}free\_\_\allowbreak{}new\_A150\_\_\allowbreak{}seed6} & lot D & campagne A/B ; régime nouveau & vintage $A\textbackslash{}times1{,}5$, sens libre \\ \textit{t=2001 A$\rightarrow$1.5 (new)} \\
\code{m4\_3live\_v2\_\_\allowbreak{}arm\_\_\allowbreak{}free\_\_\allowbreak{}new\_A150\_\_\allowbreak{}seed7} & lot D & campagne A/B ; régime nouveau & vintage $A\textbackslash{}times1{,}5$, sens libre \\ \textit{t=2001 A$\rightarrow$1.5 (new)} \\
\code{m4\_3live\_v2\_\_\allowbreak{}arm\_\_\allowbreak{}free\_\_\allowbreak{}new\_A150\_\_\allowbreak{}seed8} & lot D & campagne A/B ; régime nouveau & vintage $A\textbackslash{}times1{,}5$, sens libre \\ \textit{t=2001 A$\rightarrow$1.5 (new)} \\
\code{m4\_3live\_v2\_\_\allowbreak{}arm\_\_\allowbreak{}free\_\_\allowbreak{}new\_A150\_\_\allowbreak{}seed9} & lot D & campagne A/B ; régime nouveau & vintage $A\textbackslash{}times1{,}5$, sens libre \\ \textit{t=2001 A$\rightarrow$1.5 (new)} \\
\code{m4\_3live\_v2\_\_\allowbreak{}arm\_\_\allowbreak{}free\_\_\allowbreak{}new\_g060\_\_\allowbreak{}seed0} & lot D & campagne A/B & vintage $\textbackslash{}gamma : 0{,}5 \textbackslash{}to 0{,}6$, sens libre ; exerce le régime (c) du noyau \\ \textit{t=2001 gamma$\rightarrow$0.6 (new)} \\
\code{m4\_3live\_v2\_\_\allowbreak{}arm\_\_\allowbreak{}free\_\_\allowbreak{}new\_g060\_\_\allowbreak{}seed1} & lot D & campagne A/B & vintage $\textbackslash{}gamma : 0{,}5 \textbackslash{}to 0{,}6$, sens libre ; exerce le régime (c) du noyau \\ \textit{t=2001 gamma$\rightarrow$0.6 (new)} \\
\code{m4\_3live\_v2\_\_\allowbreak{}arm\_\_\allowbreak{}free\_\_\allowbreak{}new\_g060\_\_\allowbreak{}seed10} & lot D & campagne A/B & vintage $\textbackslash{}gamma : 0{,}5 \textbackslash{}to 0{,}6$, sens libre ; exerce le régime (c) du noyau \\ \textit{t=2001 gamma$\rightarrow$0.6 (new)} \\
\code{m4\_3live\_v2\_\_\allowbreak{}arm\_\_\allowbreak{}free\_\_\allowbreak{}new\_g060\_\_\allowbreak{}seed11} & lot D & campagne A/B & vintage $\textbackslash{}gamma : 0{,}5 \textbackslash{}to 0{,}6$, sens libre ; exerce le régime (c) du noyau \\ \textit{t=2001 gamma$\rightarrow$0.6 (new)} \\
\code{m4\_3live\_v2\_\_\allowbreak{}arm\_\_\allowbreak{}free\_\_\allowbreak{}new\_g060\_\_\allowbreak{}seed2} & lot D & campagne A/B & vintage $\textbackslash{}gamma : 0{,}5 \textbackslash{}to 0{,}6$, sens libre ; exerce le régime (c) du noyau \\ \textit{t=2001 gamma$\rightarrow$0.6 (new)} \\
\code{m4\_3live\_v2\_\_\allowbreak{}arm\_\_\allowbreak{}free\_\_\allowbreak{}new\_g060\_\_\allowbreak{}seed3} & lot D & campagne A/B & vintage $\textbackslash{}gamma : 0{,}5 \textbackslash{}to 0{,}6$, sens libre ; exerce le régime (c) du noyau \\ \textit{t=2001 gamma$\rightarrow$0.6 (new)} \\
\code{m4\_3live\_v2\_\_\allowbreak{}arm\_\_\allowbreak{}free\_\_\allowbreak{}new\_g060\_\_\allowbreak{}seed4} & lot D & campagne A/B & vintage $\textbackslash{}gamma : 0{,}5 \textbackslash{}to 0{,}6$, sens libre ; exerce le régime (c) du noyau \\ \textit{t=2001 gamma$\rightarrow$0.6 (new)} \\
\code{m4\_3live\_v2\_\_\allowbreak{}arm\_\_\allowbreak{}free\_\_\allowbreak{}new\_g060\_\_\allowbreak{}seed5} & lot D & campagne A/B & vintage $\textbackslash{}gamma : 0{,}5 \textbackslash{}to 0{,}6$, sens libre ; exerce le régime (c) du noyau \\ \textit{t=2001 gamma$\rightarrow$0.6 (new)} \\
\code{m4\_3live\_v2\_\_\allowbreak{}arm\_\_\allowbreak{}free\_\_\allowbreak{}new\_g060\_\_\allowbreak{}seed6} & lot D & campagne A/B & vintage $\textbackslash{}gamma : 0{,}5 \textbackslash{}to 0{,}6$, sens libre ; exerce le régime (c) du noyau \\ \textit{t=2001 gamma$\rightarrow$0.6 (new)} \\
\code{m4\_3live\_v2\_\_\allowbreak{}arm\_\_\allowbreak{}free\_\_\allowbreak{}new\_g060\_\_\allowbreak{}seed7} & lot D & campagne A/B & vintage $\textbackslash{}gamma : 0{,}5 \textbackslash{}to 0{,}6$, sens libre ; exerce le régime (c) du noyau \\ \textit{t=2001 gamma$\rightarrow$0.6 (new)} \\
\code{m4\_3live\_v2\_\_\allowbreak{}arm\_\_\allowbreak{}free\_\_\allowbreak{}new\_g060\_\_\allowbreak{}seed8} & lot D & campagne A/B & vintage $\textbackslash{}gamma : 0{,}5 \textbackslash{}to 0{,}6$, sens libre ; exerce le régime (c) du noyau \\ \textit{t=2001 gamma$\rightarrow$0.6 (new)} \\
\code{m4\_3live\_v2\_\_\allowbreak{}arm\_\_\allowbreak{}free\_\_\allowbreak{}new\_g060\_\_\allowbreak{}seed9} & lot D & campagne A/B & vintage $\textbackslash{}gamma : 0{,}5 \textbackslash{}to 0{,}6$, sens libre ; exerce le régime (c) du noyau \\ \textit{t=2001 gamma$\rightarrow$0.6 (new)} \\
\code{m4\_3live\_v2\_\_\allowbreak{}arm\_\_\allowbreak{}richest\_lends\_\_\allowbreak{}new\_A075\_\_\allowbreak{}seed0} & lot D & campagne A/B & même bras sous la règle v1 \\ \textit{t=2001 A$\rightarrow$0.75 (new)} \\
\code{m4\_3live\_v2\_\_\allowbreak{}arm\_\_\allowbreak{}richest\_lends\_\_\allowbreak{}new\_A075\_\_\allowbreak{}seed1} & lot D & campagne A/B & même bras sous la règle v1 \\ \textit{t=2001 A$\rightarrow$0.75 (new)} \\
\code{m4\_3live\_v2\_\_\allowbreak{}arm\_\_\allowbreak{}richest\_lends\_\_\allowbreak{}new\_A075\_\_\allowbreak{}seed10} & lot D & campagne A/B & même bras sous la règle v1 \\ \textit{t=2001 A$\rightarrow$0.75 (new)} \\
\code{m4\_3live\_v2\_\_\allowbreak{}arm\_\_\allowbreak{}richest\_lends\_\_\allowbreak{}new\_A075\_\_\allowbreak{}seed11} & lot D & campagne A/B & même bras sous la règle v1 \\ \textit{t=2001 A$\rightarrow$0.75 (new)} \\
\code{m4\_3live\_v2\_\_\allowbreak{}arm\_\_\allowbreak{}richest\_lends\_\_\allowbreak{}new\_A075\_\_\allowbreak{}seed2} & lot D & campagne A/B & même bras sous la règle v1 \\ \textit{t=2001 A$\rightarrow$0.75 (new)} \\
\code{m4\_3live\_v2\_\_\allowbreak{}arm\_\_\allowbreak{}richest\_lends\_\_\allowbreak{}new\_A075\_\_\allowbreak{}seed3} & lot D & campagne A/B & même bras sous la règle v1 \\ \textit{t=2001 A$\rightarrow$0.75 (new)} \\
\code{m4\_3live\_v2\_\_\allowbreak{}arm\_\_\allowbreak{}richest\_lends\_\_\allowbreak{}new\_A075\_\_\allowbreak{}seed4} & lot D & campagne A/B & même bras sous la règle v1 \\ \textit{t=2001 A$\rightarrow$0.75 (new)} \\
\code{m4\_3live\_v2\_\_\allowbreak{}arm\_\_\allowbreak{}richest\_lends\_\_\allowbreak{}new\_A075\_\_\allowbreak{}seed5} & lot D & campagne A/B & même bras sous la règle v1 \\ \textit{t=2001 A$\rightarrow$0.75 (new)} \\
\code{m4\_3live\_v2\_\_\allowbreak{}arm\_\_\allowbreak{}richest\_lends\_\_\allowbreak{}new\_A075\_\_\allowbreak{}seed6} & lot D & campagne A/B & même bras sous la règle v1 \\ \textit{t=2001 A$\rightarrow$0.75 (new)} \\
\code{m4\_3live\_v2\_\_\allowbreak{}arm\_\_\allowbreak{}richest\_lends\_\_\allowbreak{}new\_A075\_\_\allowbreak{}seed7} & lot D & campagne A/B & même bras sous la règle v1 \\ \textit{t=2001 A$\rightarrow$0.75 (new)} \\
\code{m4\_3live\_v2\_\_\allowbreak{}arm\_\_\allowbreak{}richest\_lends\_\_\allowbreak{}new\_A075\_\_\allowbreak{}seed8} & lot D & campagne A/B & même bras sous la règle v1 \\ \textit{t=2001 A$\rightarrow$0.75 (new)} \\
\code{m4\_3live\_v2\_\_\allowbreak{}arm\_\_\allowbreak{}richest\_lends\_\_\allowbreak{}new\_A075\_\_\allowbreak{}seed9} & lot D & campagne A/B & même bras sous la règle v1 \\ \textit{t=2001 A$\rightarrow$0.75 (new)} \\
\code{m4\_3live\_v2\_\_\allowbreak{}arm\_\_\allowbreak{}richest\_lends\_\_\allowbreak{}new\_A150\_\_\allowbreak{}seed0} & lot D & campagne A/B & même bras sous la règle v1 « la plus riche prête » \\ \textit{t=2001 A$\rightarrow$1.5 (new)} \\
\code{m4\_3live\_v2\_\_\allowbreak{}arm\_\_\allowbreak{}richest\_lends\_\_\allowbreak{}new\_A150\_\_\allowbreak{}seed1} & lot D & campagne A/B & même bras sous la règle v1 « la plus riche prête » \\ \textit{t=2001 A$\rightarrow$1.5 (new)} \\
\code{m4\_3live\_v2\_\_\allowbreak{}arm\_\_\allowbreak{}richest\_lends\_\_\allowbreak{}new\_A150\_\_\allowbreak{}seed10} & lot D & campagne A/B & même bras sous la règle v1 « la plus riche prête » \\ \textit{t=2001 A$\rightarrow$1.5 (new)} \\
\code{m4\_3live\_v2\_\_\allowbreak{}arm\_\_\allowbreak{}richest\_lends\_\_\allowbreak{}new\_A150\_\_\allowbreak{}seed11} & lot D & campagne A/B & même bras sous la règle v1 « la plus riche prête » \\ \textit{t=2001 A$\rightarrow$1.5 (new)} \\
\code{m4\_3live\_v2\_\_\allowbreak{}arm\_\_\allowbreak{}richest\_lends\_\_\allowbreak{}new\_A150\_\_\allowbreak{}seed2} & lot D & campagne A/B & même bras sous la règle v1 « la plus riche prête » \\ \textit{t=2001 A$\rightarrow$1.5 (new)} \\
\code{m4\_3live\_v2\_\_\allowbreak{}arm\_\_\allowbreak{}richest\_lends\_\_\allowbreak{}new\_A150\_\_\allowbreak{}seed3} & lot D & campagne A/B & même bras sous la règle v1 « la plus riche prête » \\ \textit{t=2001 A$\rightarrow$1.5 (new)} \\
\code{m4\_3live\_v2\_\_\allowbreak{}arm\_\_\allowbreak{}richest\_lends\_\_\allowbreak{}new\_A150\_\_\allowbreak{}seed4} & lot D & campagne A/B & même bras sous la règle v1 « la plus riche prête » \\ \textit{t=2001 A$\rightarrow$1.5 (new)} \\
\code{m4\_3live\_v2\_\_\allowbreak{}arm\_\_\allowbreak{}richest\_lends\_\_\allowbreak{}new\_A150\_\_\allowbreak{}seed5} & lot D & campagne A/B & même bras sous la règle v1 « la plus riche prête » \\ \textit{t=2001 A$\rightarrow$1.5 (new)} \\
\code{m4\_3live\_v2\_\_\allowbreak{}arm\_\_\allowbreak{}richest\_lends\_\_\allowbreak{}new\_A150\_\_\allowbreak{}seed6} & lot D & campagne A/B & même bras sous la règle v1 « la plus riche prête » \\ \textit{t=2001 A$\rightarrow$1.5 (new)} \\
\code{m4\_3live\_v2\_\_\allowbreak{}arm\_\_\allowbreak{}richest\_lends\_\_\allowbreak{}new\_A150\_\_\allowbreak{}seed7} & lot D & campagne A/B & même bras sous la règle v1 « la plus riche prête » \\ \textit{t=2001 A$\rightarrow$1.5 (new)} \\
\code{m4\_3live\_v2\_\_\allowbreak{}arm\_\_\allowbreak{}richest\_lends\_\_\allowbreak{}new\_A150\_\_\allowbreak{}seed8} & lot D & campagne A/B & même bras sous la règle v1 « la plus riche prête » \\ \textit{t=2001 A$\rightarrow$1.5 (new)} \\
\code{m4\_3live\_v2\_\_\allowbreak{}arm\_\_\allowbreak{}richest\_lends\_\_\allowbreak{}new\_A150\_\_\allowbreak{}seed9} & lot D & campagne A/B & même bras sous la règle v1 « la plus riche prête » \\ \textit{t=2001 A$\rightarrow$1.5 (new)} \\
\code{m4\_3live\_v2\_\_\allowbreak{}arm\_\_\allowbreak{}richest\_lends\_\_\allowbreak{}new\_g060\_\_\allowbreak{}seed0} & lot D & campagne A/B & même bras sous la règle v1 \\ \textit{t=2001 gamma$\rightarrow$0.6 (new)} \\
\code{m4\_3live\_v2\_\_\allowbreak{}arm\_\_\allowbreak{}richest\_lends\_\_\allowbreak{}new\_g060\_\_\allowbreak{}seed1} & lot D & campagne A/B & même bras sous la règle v1 \\ \textit{t=2001 gamma$\rightarrow$0.6 (new)} \\
\code{m4\_3live\_v2\_\_\allowbreak{}arm\_\_\allowbreak{}richest\_lends\_\_\allowbreak{}new\_g060\_\_\allowbreak{}seed10} & lot D & campagne A/B & même bras sous la règle v1 \\ \textit{t=2001 gamma$\rightarrow$0.6 (new)} \\
\code{m4\_3live\_v2\_\_\allowbreak{}arm\_\_\allowbreak{}richest\_lends\_\_\allowbreak{}new\_g060\_\_\allowbreak{}seed11} & lot D & campagne A/B & même bras sous la règle v1 \\ \textit{t=2001 gamma$\rightarrow$0.6 (new)} \\
\code{m4\_3live\_v2\_\_\allowbreak{}arm\_\_\allowbreak{}richest\_lends\_\_\allowbreak{}new\_g060\_\_\allowbreak{}seed2} & lot D & campagne A/B & même bras sous la règle v1 \\ \textit{t=2001 gamma$\rightarrow$0.6 (new)} \\
\code{m4\_3live\_v2\_\_\allowbreak{}arm\_\_\allowbreak{}richest\_lends\_\_\allowbreak{}new\_g060\_\_\allowbreak{}seed3} & lot D & campagne A/B & même bras sous la règle v1 \\ \textit{t=2001 gamma$\rightarrow$0.6 (new)} \\
\code{m4\_3live\_v2\_\_\allowbreak{}arm\_\_\allowbreak{}richest\_lends\_\_\allowbreak{}new\_g060\_\_\allowbreak{}seed4} & lot D & campagne A/B & même bras sous la règle v1 \\ \textit{t=2001 gamma$\rightarrow$0.6 (new)} \\
\code{m4\_3live\_v2\_\_\allowbreak{}arm\_\_\allowbreak{}richest\_lends\_\_\allowbreak{}new\_g060\_\_\allowbreak{}seed5} & lot D & campagne A/B & même bras sous la règle v1 \\ \textit{t=2001 gamma$\rightarrow$0.6 (new)} \\
\code{m4\_3live\_v2\_\_\allowbreak{}arm\_\_\allowbreak{}richest\_lends\_\_\allowbreak{}new\_g060\_\_\allowbreak{}seed6} & lot D & campagne A/B & même bras sous la règle v1 \\ \textit{t=2001 gamma$\rightarrow$0.6 (new)} \\
\code{m4\_3live\_v2\_\_\allowbreak{}arm\_\_\allowbreak{}richest\_lends\_\_\allowbreak{}new\_g060\_\_\allowbreak{}seed7} & lot D & campagne A/B & même bras sous la règle v1 \\ \textit{t=2001 gamma$\rightarrow$0.6 (new)} \\
\code{m4\_3live\_v2\_\_\allowbreak{}arm\_\_\allowbreak{}richest\_lends\_\_\allowbreak{}new\_g060\_\_\allowbreak{}seed8} & lot D & campagne A/B & même bras sous la règle v1 \\ \textit{t=2001 gamma$\rightarrow$0.6 (new)} \\
\code{m4\_3live\_v2\_\_\allowbreak{}arm\_\_\allowbreak{}richest\_lends\_\_\allowbreak{}new\_g060\_\_\allowbreak{}seed9} & lot D & campagne A/B & même bras sous la règle v1 \\ \textit{t=2001 gamma$\rightarrow$0.6 (new)} \\
\addlinespace
\code{m4\_3live\_v2\_\_\allowbreak{}burn\_\_\allowbreak{}seed0} & amorçage & protocole ; source des snapshots à $t_0 = 2000$ de tous les bras & amorçage partagé 0 $\rightarrow$ 2000, homogène, commun aux deux règles de sens \\
\code{m4\_3live\_v2\_\_\allowbreak{}burn\_\_\allowbreak{}seed1} & amorçage & protocole ; source des snapshots à $t_0 = 2000$ de tous les bras & amorçage partagé 0 $\rightarrow$ 2000, homogène, commun aux deux règles de sens \\
\code{m4\_3live\_v2\_\_\allowbreak{}burn\_\_\allowbreak{}seed10} & amorçage & protocole ; source des snapshots à $t_0 = 2000$ de tous les bras & amorçage partagé 0 $\rightarrow$ 2000, homogène, commun aux deux règles de sens \\
\code{m4\_3live\_v2\_\_\allowbreak{}burn\_\_\allowbreak{}seed11} & amorçage & protocole ; source des snapshots à $t_0 = 2000$ de tous les bras & amorçage partagé 0 $\rightarrow$ 2000, homogène, commun aux deux règles de sens \\
\code{m4\_3live\_v2\_\_\allowbreak{}burn\_\_\allowbreak{}seed2} & amorçage & protocole ; source des snapshots à $t_0 = 2000$ de tous les bras & amorçage partagé 0 $\rightarrow$ 2000, homogène, commun aux deux règles de sens \\
\code{m4\_3live\_v2\_\_\allowbreak{}burn\_\_\allowbreak{}seed3} & amorçage & protocole ; source des snapshots à $t_0 = 2000$ de tous les bras & amorçage partagé 0 $\rightarrow$ 2000, homogène, commun aux deux règles de sens \\
\code{m4\_3live\_v2\_\_\allowbreak{}burn\_\_\allowbreak{}seed4} & amorçage & protocole ; source des snapshots à $t_0 = 2000$ de tous les bras & amorçage partagé 0 $\rightarrow$ 2000, homogène, commun aux deux règles de sens \\
\code{m4\_3live\_v2\_\_\allowbreak{}burn\_\_\allowbreak{}seed5} & amorçage & protocole ; source des snapshots à $t_0 = 2000$ de tous les bras & amorçage partagé 0 $\rightarrow$ 2000, homogène, commun aux deux règles de sens \\
\code{m4\_3live\_v2\_\_\allowbreak{}burn\_\_\allowbreak{}seed6} & amorçage & protocole ; source des snapshots à $t_0 = 2000$ de tous les bras & amorçage partagé 0 $\rightarrow$ 2000, homogène, commun aux deux règles de sens \\
\code{m4\_3live\_v2\_\_\allowbreak{}burn\_\_\allowbreak{}seed7} & amorçage & protocole ; source des snapshots à $t_0 = 2000$ de tous les bras & amorçage partagé 0 $\rightarrow$ 2000, homogène, commun aux deux règles de sens \\
\code{m4\_3live\_v2\_\_\allowbreak{}burn\_\_\allowbreak{}seed8} & amorçage & protocole ; source des snapshots à $t_0 = 2000$ de tous les bras & amorçage partagé 0 $\rightarrow$ 2000, homogène, commun aux deux règles de sens \\
\code{m4\_3live\_v2\_\_\allowbreak{}burn\_\_\allowbreak{}seed9} & amorçage & protocole ; source des snapshots à $t_0 = 2000$ de tous les bras & amorçage partagé 0 $\rightarrow$ 2000, homogène, commun aux deux règles de sens \\
\addlinespace
\code{m4\_3live\_v2\_\_\allowbreak{}phase\_\_\allowbreak{}deprec\_first\_\_\allowbreak{}seed0} & lot E & ordre des phases & dépréciation avant service des intérêts, apparié au contrôle \\
\code{m4\_3live\_v2\_\_\allowbreak{}phase\_\_\allowbreak{}deprec\_first\_\_\allowbreak{}seed1} & lot E & ordre des phases & dépréciation avant service des intérêts, apparié au contrôle \\
\code{m4\_3live\_v2\_\_\allowbreak{}phase\_\_\allowbreak{}deprec\_first\_\_\allowbreak{}seed10} & lot E & ordre des phases & dépréciation avant service des intérêts, apparié au contrôle \\
\code{m4\_3live\_v2\_\_\allowbreak{}phase\_\_\allowbreak{}deprec\_first\_\_\allowbreak{}seed11} & lot E & ordre des phases & dépréciation avant service des intérêts, apparié au contrôle \\
\code{m4\_3live\_v2\_\_\allowbreak{}phase\_\_\allowbreak{}deprec\_first\_\_\allowbreak{}seed2} & lot E & ordre des phases & dépréciation avant service des intérêts, apparié au contrôle \\
\code{m4\_3live\_v2\_\_\allowbreak{}phase\_\_\allowbreak{}deprec\_first\_\_\allowbreak{}seed3} & lot E & ordre des phases & dépréciation avant service des intérêts, apparié au contrôle \\
\code{m4\_3live\_v2\_\_\allowbreak{}phase\_\_\allowbreak{}deprec\_first\_\_\allowbreak{}seed4} & lot E & ordre des phases & dépréciation avant service des intérêts, apparié au contrôle \\
\code{m4\_3live\_v2\_\_\allowbreak{}phase\_\_\allowbreak{}deprec\_first\_\_\allowbreak{}seed5} & lot E & ordre des phases & dépréciation avant service des intérêts, apparié au contrôle \\
\code{m4\_3live\_v2\_\_\allowbreak{}phase\_\_\allowbreak{}deprec\_first\_\_\allowbreak{}seed6} & lot E & ordre des phases & dépréciation avant service des intérêts, apparié au contrôle \\
\code{m4\_3live\_v2\_\_\allowbreak{}phase\_\_\allowbreak{}deprec\_first\_\_\allowbreak{}seed7} & lot E & ordre des phases & dépréciation avant service des intérêts, apparié au contrôle \\
\code{m4\_3live\_v2\_\_\allowbreak{}phase\_\_\allowbreak{}deprec\_first\_\_\allowbreak{}seed8} & lot E & ordre des phases & dépréciation avant service des intérêts, apparié au contrôle \\
\code{m4\_3live\_v2\_\_\allowbreak{}phase\_\_\allowbreak{}deprec\_first\_\_\allowbreak{}seed9} & lot E & ordre des phases & dépréciation avant service des intérêts, apparié au contrôle \\
\addlinespace
\code{m4\_3live\_v2\_\_\allowbreak{}rotation\_\_\allowbreak{}K0\_100\_\_\allowbreak{}seed0} & lot F & rotation du crédit (décomposition et fermeture) & levier $K\_0$ (capital de naissance) seul ($= 100$) \\
\code{m4\_3live\_v2\_\_\allowbreak{}rotation\_\_\allowbreak{}K0\_100\_\_\allowbreak{}seed1} & lot F & rotation du crédit (décomposition et fermeture) & levier $K\_0$ (capital de naissance) seul ($= 100$) \\
\code{m4\_3live\_v2\_\_\allowbreak{}rotation\_\_\allowbreak{}K0\_100\_\_\allowbreak{}seed2} & lot F & rotation du crédit (décomposition et fermeture) & levier $K\_0$ (capital de naissance) seul ($= 100$) \\
\code{m4\_3live\_v2\_\_\allowbreak{}rotation\_\_\allowbreak{}K0\_6.25\_\_\allowbreak{}seed0} & lot F & rotation du crédit (décomposition et fermeture) & levier $K\_0$ (capital de naissance) seul ($= 6,25$) \\
\code{m4\_3live\_v2\_\_\allowbreak{}rotation\_\_\allowbreak{}K0\_6.25\_\_\allowbreak{}seed1} & lot F & rotation du crédit (décomposition et fermeture) & levier $K\_0$ (capital de naissance) seul ($= 6,25$) \\
\code{m4\_3live\_v2\_\_\allowbreak{}rotation\_\_\allowbreak{}K0\_6.25\_\_\allowbreak{}seed2} & lot F & rotation du crédit (décomposition et fermeture) & levier $K\_0$ (capital de naissance) seul ($= 6,25$) \\
\code{m4\_3live\_v2\_\_\allowbreak{}rotation\_\_\allowbreak{}base\_\_\allowbreak{}seed0} & lot F & rotation du crédit (décomposition et fermeture) & régime de référence, avec le Gini du capital instrumenté \\
\code{m4\_3live\_v2\_\_\allowbreak{}rotation\_\_\allowbreak{}base\_\_\allowbreak{}seed1} & lot F & rotation du crédit (décomposition et fermeture) & régime de référence, avec le Gini du capital instrumenté \\
\code{m4\_3live\_v2\_\_\allowbreak{}rotation\_\_\allowbreak{}base\_\_\allowbreak{}seed2} & lot F & rotation du crédit (décomposition et fermeture) & régime de référence, avec le Gini du capital instrumenté \\
\code{m4\_3live\_v2\_\_\allowbreak{}rotation\_\_\allowbreak{}delta\_0.005\_\_\allowbreak{}seed0} & lot F & rotation du crédit (décomposition et fermeture) & levier $\textbackslash{}delta$ (dépréciation) seul ($= 0,005$) \\
\code{m4\_3live\_v2\_\_\allowbreak{}rotation\_\_\allowbreak{}delta\_0.005\_\_\allowbreak{}seed1} & lot F & rotation du crédit (décomposition et fermeture) & levier $\textbackslash{}delta$ (dépréciation) seul ($= 0,005$) \\
\code{m4\_3live\_v2\_\_\allowbreak{}rotation\_\_\allowbreak{}delta\_0.005\_\_\allowbreak{}seed2} & lot F & rotation du crédit (décomposition et fermeture) & levier $\textbackslash{}delta$ (dépréciation) seul ($= 0,005$) \\
\code{m4\_3live\_v2\_\_\allowbreak{}rotation\_\_\allowbreak{}delta\_0.02\_\_\allowbreak{}seed0} & lot F & rotation du crédit (décomposition et fermeture) & levier $\textbackslash{}delta$ (dépréciation) seul ($= 0,02$) \\
\code{m4\_3live\_v2\_\_\allowbreak{}rotation\_\_\allowbreak{}delta\_0.02\_\_\allowbreak{}seed1} & lot F & rotation du crédit (décomposition et fermeture) & levier $\textbackslash{}delta$ (dépréciation) seul ($= 0,02$) \\
\code{m4\_3live\_v2\_\_\allowbreak{}rotation\_\_\allowbreak{}delta\_0.02\_\_\allowbreak{}seed2} & lot F & rotation du crédit (décomposition et fermeture) & levier $\textbackslash{}delta$ (dépréciation) seul ($= 0,02$) \\
\code{m4\_3live\_v2\_\_\allowbreak{}rotation\_\_\allowbreak{}lam\_15\_\_\allowbreak{}seed0} & lot F & rotation du crédit (décomposition et fermeture) & levier $\textbackslash{}lambda$ (taux de naissance) seul ($= 15$) \\
\code{m4\_3live\_v2\_\_\allowbreak{}rotation\_\_\allowbreak{}lam\_15\_\_\allowbreak{}seed1} & lot F & rotation du crédit (décomposition et fermeture) & levier $\textbackslash{}lambda$ (taux de naissance) seul ($= 15$) \\
\code{m4\_3live\_v2\_\_\allowbreak{}rotation\_\_\allowbreak{}lam\_15\_\_\allowbreak{}seed2} & lot F & rotation du crédit (décomposition et fermeture) & levier $\textbackslash{}lambda$ (taux de naissance) seul ($= 15$) \\
\code{m4\_3live\_v2\_\_\allowbreak{}rotation\_\_\allowbreak{}lam\_60\_\_\allowbreak{}seed0} & lot F & rotation du crédit (décomposition et fermeture) & levier $\textbackslash{}lambda$ (taux de naissance) seul ($= 60$) \\
\code{m4\_3live\_v2\_\_\allowbreak{}rotation\_\_\allowbreak{}lam\_60\_\_\allowbreak{}seed1} & lot F & rotation du crédit (décomposition et fermeture) & levier $\textbackslash{}lambda$ (taux de naissance) seul ($= 60$) \\
\code{m4\_3live\_v2\_\_\allowbreak{}rotation\_\_\allowbreak{}lam\_60\_\_\allowbreak{}seed2} & lot F & rotation du crédit (décomposition et fermeture) & levier $\textbackslash{}lambda$ (taux de naissance) seul ($= 60$) \\
\code{m4\_3live\_v2\_\_\allowbreak{}rotation\_\_\allowbreak{}rho\_0.5\_\_\allowbreak{}seed0} & lot F & rotation du crédit (décomposition et fermeture) & levier $\textbackslash{}rho$ (taux d'appariement) seul ($= 0,5$) \\
\code{m4\_3live\_v2\_\_\allowbreak{}rotation\_\_\allowbreak{}rho\_0.5\_\_\allowbreak{}seed1} & lot F & rotation du crédit (décomposition et fermeture) & levier $\textbackslash{}rho$ (taux d'appariement) seul ($= 0,5$) \\
\code{m4\_3live\_v2\_\_\allowbreak{}rotation\_\_\allowbreak{}rho\_0.5\_\_\allowbreak{}seed2} & lot F & rotation du crédit (décomposition et fermeture) & levier $\textbackslash{}rho$ (taux d'appariement) seul ($= 0,5$) \\
\code{m4\_3live\_v2\_\_\allowbreak{}rotation\_\_\allowbreak{}rho\_2\_\_\allowbreak{}seed0} & lot F & rotation du crédit (décomposition et fermeture) & levier $\textbackslash{}rho$ (taux d'appariement) seul ($= 2$) \\
\code{m4\_3live\_v2\_\_\allowbreak{}rotation\_\_\allowbreak{}rho\_2\_\_\allowbreak{}seed1} & lot F & rotation du crédit (décomposition et fermeture) & levier $\textbackslash{}rho$ (taux d'appariement) seul ($= 2$) \\
\code{m4\_3live\_v2\_\_\allowbreak{}rotation\_\_\allowbreak{}rho\_2\_\_\allowbreak{}seed2} & lot F & rotation du crédit (décomposition et fermeture) & levier $\textbackslash{}rho$ (taux d'appariement) seul ($= 2$) \\
\code{m4\_3live\_v2\_\_\allowbreak{}rotation\_\_\allowbreak{}sigma\_0.005\_\_\allowbreak{}seed0} & lot F & rotation du crédit (décomposition et fermeture) & levier $\textbackslash{}sigma$ (choc multiplicatif) seul ($= 0,005$) \\
\code{m4\_3live\_v2\_\_\allowbreak{}rotation\_\_\allowbreak{}sigma\_0.005\_\_\allowbreak{}seed1} & lot F & rotation du crédit (décomposition et fermeture) & levier $\textbackslash{}sigma$ (choc multiplicatif) seul ($= 0,005$) \\
\code{m4\_3live\_v2\_\_\allowbreak{}rotation\_\_\allowbreak{}sigma\_0.005\_\_\allowbreak{}seed2} & lot F & rotation du crédit (décomposition et fermeture) & levier $\textbackslash{}sigma$ (choc multiplicatif) seul ($= 0,005$) \\
\code{m4\_3live\_v2\_\_\allowbreak{}rotation\_\_\allowbreak{}sigma\_0.02\_\_\allowbreak{}seed0} & lot F & rotation du crédit (décomposition et fermeture) & levier $\textbackslash{}sigma$ (choc multiplicatif) seul ($= 0,02$) \\
\code{m4\_3live\_v2\_\_\allowbreak{}rotation\_\_\allowbreak{}sigma\_0.02\_\_\allowbreak{}seed1} & lot F & rotation du crédit (décomposition et fermeture) & levier $\textbackslash{}sigma$ (choc multiplicatif) seul ($= 0,02$) \\
\code{m4\_3live\_v2\_\_\allowbreak{}rotation\_\_\allowbreak{}sigma\_0.02\_\_\allowbreak{}seed2} & lot F & rotation du crédit (décomposition et fermeture) & levier $\textbackslash{}sigma$ (choc multiplicatif) seul ($= 0,02$) \\
\addlinespace
\code{m4\_3\_\_\allowbreak{}d1\_\_\allowbreak{}baseline\_\_\allowbreak{}seed0} & référence M4.3 & conception (parité bit à bit sur 8000 pas) ; lots A et B & run M4.3 stocké servant de référence de parité \\
\code{m4\_3\_\_\allowbreak{}d1\_\_\allowbreak{}baseline\_\_\allowbreak{}seed1} & référence M4.3 & protocole (justification de $t_0$) & t\_converge\_int\_in = 948, lu dans son analysis.json \\
\code{m4\_3\_\_\allowbreak{}d1\_\_\allowbreak{}baseline\_\_\allowbreak{}seed2} & référence M4.3 & protocole (justification de $t_0$) & t\_converge\_int\_in = 841, lu dans son analysis.json \\
\bottomrule
\end{longtable}
}

\noindent Tous ces runs sont ouvrables dans \code{simulation\_lab} (\code{python3 -m simulation\_lab.cli gui}, page « Résultats », lignée \code{m4\_3live\_v2\_credit\_soc}), avec leur \code{series.csv}, leur \code{tech\_series.csv}, leurs \code{tension.csv} et \code{tension\_agg.csv}, leur journal d'interventions et une figure \code{figures/macro\_overview.png}. Le tableau complet, avec les paramètres et le statut de chaque run, est dans \code{results/analysis/traceability.csv}.


\end{document}
